\chapter{Greedy Algorithms}
\chapterepigraph{A greedy algorithm never looks back. At every step it makes the choice that
looks best at that instant, and trusts -- or, better, proves -- that this sequence
of locally best choices assembles itself into a globally best solution.}

\section{What Is a Greedy Algorithm?}

A \textbf{greedy algorithm} builds up a solution piece by piece, and at every step
it commits, without revision, to whichever available choice looks most
favourable \emph{right now}. It never reconsiders an earlier decision, and it
never explores alternative branches the way recursion or dynamic programming
would. This makes greedy algorithms extremely fast and extremely simple to
code -- usually a single pass over the input, sometimes after a sort -- but it
also means that greedy is dangerous to apply carelessly. A locally best choice
is only safe if it can be \emph{proved} never to block a globally best
solution. When that proof exists, greedy is the right tool, and it is almost
always the most elegant one. When that proof does not exist, a greedy-looking
shortcut will quietly produce wrong answers on some hidden test case, while
looking completely correct on the examples you tried by hand.

\begin{ideabox}[When Is Greedy Safe to Use?]
A problem is a good candidate for a greedy algorithm when it has two
properties.
\begin{enumerate}
  \item \textbf{Greedy choice property.} There is a way to pick, at each step,
        a single locally optimal choice such that some globally optimal
        solution is consistent with that choice -- i.e.\ making the greedy
        choice never rules out the best possible final answer.
  \item \textbf{Optimal substructure.} Once the greedy choice is made, what
        remains is a smaller version of the same problem, and an optimal
        solution to the whole problem is built from the greedy choice plus an
        optimal solution to this smaller remaining problem.
\end{enumerate}
Contrast this with dynamic programming, which is needed precisely when the
greedy choice property \emph{fails} -- when the locally best choice at a step
can depend on, or conflict with, choices made elsewhere, forcing us to compare
many overlapping subproblems instead of committing to one path.
\end{ideabox}

\subsection*{How to Recognise a Greedy Problem from Its Wording}

Across the problems in this chapter, certain words and phrases recur whenever
greedy is the intended technique. Training yourself to notice them is the
fastest way to identify the right pattern within seconds of reading a problem
statement.

\begin{patternbox}[Recurring Greedy Keywords]
\begin{itemize}
  \item \textbf{``Maximum number of non-overlapping ...''} / \textbf{``minimum
        number of ... to remove''} -- classic interval scheduling; sort by end
        time or by a similar criterion and sweep once.
  \item \textbf{``Assign / allocate / distribute to satisfy as many as
        possible''} -- two-pointer style matching after sorting both sides
        (e.g.\ Assign Cookies).
  \item \textbf{``In order'' / ``queue'' / ``one at a time, immediately''} --
        signals that decisions must be made online, in a fixed order, with no
        ability to rearrange (e.g.\ Lemonade Change).
  \item \textbf{``Maximize profit/value subject to a capacity/weight/time
        limit'' with divisible items} -- Fractional Knapsack: take by best
        ratio of value to cost.
  \item \textbf{``Reach the last index'' / ``minimum jumps''} -- range/horizon
        extension greedy (Jump Game family).
  \item \textbf{``Schedule jobs with deadlines and profits''} -- Job
        Sequencing: sort by profit, place as late as possible before the
        deadline.
  \item \textbf{``Neighbouring elements must be strictly greater/smaller''} --
        local comparison greedy done in two directional passes (Candy).
\end{itemize}
Whenever you see sorting followed by a single linear sweep with no
backtracking, and the problem asks for an optimum count, profit, or
feasibility, greedy should be the first technique you test -- and you should
immediately try to articulate \emph{why} the locally best choice cannot hurt
the global answer, since that argument is the real content of a greedy
solution, not just the code.
\end{patternbox}

Each section in this chapter follows the same five-part structure: first we
restate the problem precisely and identify the keywords that flag it as
greedy; then we reason our way, case by case, to the greedy choice and argue
informally why it is safe; then we write the algorithm as a numbered
procedure and as formal pseudocode; then we trace the algorithm by hand on a
concrete example; and finally we give complete, compilable C++ code together
with a careful time and space complexity analysis.

\newpage

\section{Assign Cookies}
\label{sec:assign-cookies}

\problemstatement{We are given two arrays. The array $g$ of length $n$ holds
the \emph{greed factor} of each child: child $i$ will be content only if it
receives a cookie of size at least $g[i]$. The array $s$ of length $m$ holds
the \emph{size} of each available cookie. Each cookie may be given to at most
one child, and each child may receive at most one cookie. We must hand out
cookies so as to \textbf{maximize the number of content children}, and return
that maximum count.}

\begin{patternbox}
The phrase \emph{``maximize the number of children you can satisfy''} together
with two independent arrays that must be \emph{matched} against each other is
the signature of a two-pointer greedy assignment problem. Whenever a problem
gives you two collections, a one-to-one matching rule of the form ``resource
$j$ can serve request $i$ if resource $j$ is large enough'', and asks for the
maximum number of satisfied requests, the standard reflex is: \emph{sort both
arrays, then match the smallest unsatisfied request with the smallest
resource capable of satisfying it.} No recursion, no DP table, and no
exhaustive matching is needed -- sorting already throws away every irrelevant
ordering and leaves only a single greedy sweep.
\end{patternbox}

\subsection*{Understanding the problem carefully}

Strip away the cookie/child story and the problem is really about two
multisets of numbers, $g$ and $s$. We are allowed to pair an element of $g$
with an element of $s$ exactly when
\[
  s[j] \;\ge\; g[i],
\]
and each element of either array can be used in at most one pair. We want the
largest possible number of valid pairs.

Notice two things immediately.

\begin{enumerate}
  \item The actual \emph{order} in which children or cookies are listed in
        the input is irrelevant to the answer -- only the multiset of greed
        factors and the multiset of cookie sizes matter. This is a strong hint
        that we are free to reorder (sort) the input without losing anything,
        which is exactly the kind of freedom a greedy algorithm wants to
        exploit.
  \item A large cookie can always satisfy any child that a smaller cookie
        could satisfy. So if we are ever tempted to ``save'' a large cookie
        for later while a smaller cookie sitting right next to it could have
        done the same job, we are never making things better, and sometimes
        we are making things strictly worse, because the large cookie might
        have been the \emph{only} cookie capable of satisfying some greedier
        child later on.
\end{enumerate}

\subsection*{Why sorting is the first move}

\begin{ideabox}[Greedy Choice]
Sort $g$ and $s$ in increasing order. Walk through the children from the
least greedy to the most greedy. For the current least-greedy unsatisfied
child, give them the \emph{smallest cookie that is still large enough} to
satisfy them. If no remaining cookie is large enough, this child -- and every
child after them, since they are all at least as greedy -- cannot be the next
one satisfied with the cookies left, so we move on without giving up a
cookie unnecessarily.
\end{ideabox}

Why does always matching the least-greedy child with the smallest sufficient
cookie never hurt? Suppose, for contradiction, that some optimal assignment
does \emph{not} give the least-greedy child $c$ the smallest sufficient
cookie $k$ available, but instead gives $c$ some larger sufficient cookie
$k'$, while $k$ is given to a different child $c'$ (or not used at all). Since
$c$ is the least greedy child being considered, and $k$ is already known to
satisfy $c$ (it is sufficient), and $k$ is the smallest such cookie, swapping
$k$ and $k'$ between $c$ and $c'$ keeps $c$ satisfied (since $k$ is still
sufficient for $c$ by definition) and does not break $c'$'s satisfaction
either, because we only handed $c'$ a \emph{larger} cookie $k'$ than before --
and a larger cookie satisfies anyone a smaller one would. So the swap never
decreases the number of satisfied children. This is a standard
\emph{exchange argument}: any optimal solution can be rearranged, one swap at
a time, into the greedy solution without ever lowering the count of satisfied
children. That is precisely what makes the greedy choice safe.

\subsection*{Why two pointers are enough}

Once both arrays are sorted, we never need to look backwards. We keep one
pointer $i$ over children (starting at the least greedy) and one pointer $j$
over cookies (starting at the smallest cookie).

\begin{itemize}
  \item If the current cookie $s[j]$ is large enough for the current child
        $g[i]$, this is the cheapest cookie that can possibly satisfy this
        child (every cookie before $j$ was too small, by sorted order), so we
        commit: give it to the child, and advance both $i$ and $j$.
  \item If $s[j]$ is too small for $g[i]$, this cookie is too small for
        \emph{every} remaining child too (they are all at least as greedy as
        $g[i]$), so it can never be used -- discard it by advancing $j$ alone.
\end{itemize}

Each pointer only ever moves forward, and each comparison either satisfies a
child or permanently discards a cookie. There is no benefit to revisiting a
discarded cookie or an already-satisfied child, so a single linear sweep after
sorting is provably sufficient -- no recursion or dynamic programming is
needed because there is no overlapping subproblem to memoize; every decision
is final the moment it is made.

\subsection*{Algorithm}

\begin{enumerate}
  \item Sort $g$ in increasing order; sort $s$ in increasing order.
  \item Set $i \gets 0$ (pointer into $g$), $j \gets 0$ (pointer into $s$),
        $\textit{count} \gets 0$.
  \item While $i < n$ and $j < m$:
  \begin{enumerate}
    \item If $s[j] \ge g[i]$: this cookie satisfies this child. Increment
          \textit{count}, and advance both $i \gets i+1$ and $j \gets j+1$.
    \item Otherwise: this cookie is too small for every remaining child.
          Advance $j \gets j+1$ only.
  \end{enumerate}
  \item Return \textit{count}.
\end{enumerate}

\begin{algorithm}[H]
\caption{Assign Cookies}
\begin{algorithmic}[1]
\Procedure{FindContentChildren}{$g, s$}
  \State sort $g$ in increasing order
  \State sort $s$ in increasing order
  \State $i \gets 0$, $j \gets 0$, $\textit{count} \gets 0$
  \While{$i < \text{length}(g)$ \textbf{and} $j < \text{length}(s)$}
    \If{$s[j] \ge g[i]$}
      \State $\textit{count} \gets \textit{count} + 1$
      \State $i \gets i + 1$
      \State $j \gets j + 1$
    \Else
      \State $j \gets j + 1$
    \EndIf
  \EndWhile
  \State \Return \textit{count}
\EndProcedure
\end{algorithmic}
\end{algorithm}

\subsection*{Dry run}

\dryrun{Take $g = [1, 2, 3]$ and $s = [1, 1]$.}

Both arrays are already sorted. We initialize $i = 0$, $j = 0$,
$\textit{count} = 0$.

\begin{itemize}
  \item $i=0, j=0$: compare $s[0]=1$ with $g[0]=1$. Since $1 \ge 1$, the
        child is satisfied. $\textit{count} = 1$, $i \to 1$, $j \to 1$.
  \item $i=1, j=1$: compare $s[1]=1$ with $g[1]=2$. Since $1 < 2$, this
        cookie is too small. $j \to 2$.
  \item $j = 2$ equals the length of $s$, so the loop ends.
\end{itemize}

The final answer is $\textit{count} = 1$. Indeed, with cookies of size
$[1,1]$ we can satisfy only the child with the smallest greed factor $1$; the
children needing $2$ and $3$ can never be satisfied by a cookie of size $1$.

Now take a second example: $g = [1, 2]$ and $s = [1, 2, 3]$.

\begin{itemize}
  \item $i=0, j=0$: $s[0]=1 \ge g[0]=1$. Satisfied. $\textit{count}=1$,
        $i \to 1$, $j \to 1$.
  \item $i=1, j=1$: $s[1]=2 \ge g[1]=2$. Satisfied. $\textit{count}=2$,
        $i \to 2$, $j \to 2$.
  \item $i = 2$ equals the length of $g$, so the loop ends.
\end{itemize}

The final answer is $\textit{count} = 2$: every child is satisfied, and the
leftover cookie of size $3$ is simply not needed.

\subsection*{C++ implementation}

\begin{lstlisting}
#include <bits/stdc++.h>
using namespace std;

int findContentChildren(vector<int>& g, vector<int>& s) {
    sort(g.begin(), g.end());
    sort(s.begin(), s.end());

    int i = 0, j = 0, count = 0;
    int n = g.size(), m = s.size();

    while (i < n && j < m) {
        if (s[j] >= g[i]) {
            count++;
            i++;
            j++;
        } else {
            j++;
        }
    }
    return count;
}

int main() {
    vector<int> g = {1, 2, 3};
    vector<int> s = {1, 1};

    cout << "Number of content children: "
         << findContentChildren(g, s) << endl;
    return 0;
}
\end{lstlisting}

\begin{complexitybox}
\textbf{Time Complexity.} Sorting $g$ takes $O(n \log n)$ and sorting $s$
takes $O(m \log m)$. The subsequent two-pointer sweep visits each index of
$g$ and $s$ at most once, so it costs $O(n + m)$. The sorting step dominates,
giving an overall time complexity of
\[
  O(n \log n + m \log m).
\]
\textbf{Space Complexity.} Sorting can be done in place, and we use only the
two pointers $i, j$ and the counter as extra variables, so the auxiliary
space complexity is $O(1)$ (ignoring the in-place sort's recursion stack,
which is $O(\log n)$ in the worst case for typical \texttt{std::sort}
implementations).
\end{complexitybox}

\begin{takeawaybox}
When a problem asks you to \emph{match} two sorted resources against each
other to maximize the count of satisfied requirements, sort both sides and
sweep with two pointers, always trying to satisfy the smallest unsatisfied
requirement with the cheapest resource capable of doing so. This pattern --
\emph{sort, then greedily match smallest-feasible-to-smallest-need} -- recurs
throughout greedy algorithms and is worth recognizing on sight.
\end{takeawaybox}

\section{Fractional Knapsack Problem}
\label{sec:fractional-knapsack}

\problemstatement{We are given $n$ items, where the $i$-th item has a value
$\textit{val}[i]$ and a weight $\textit{wt}[i]$, and a knapsack of capacity
$W$. Unlike the classical $0/1$ knapsack, here we are allowed to take any
\emph{fraction} of an item -- if we take a fraction $x \in [0,1]$ of item $i$,
we gain $x \cdot \textit{val}[i]$ value while consuming $x \cdot
\textit{wt}[i]$ weight. We must choose fractions for each item, subject to
the total weight taken not exceeding $W$, so as to \textbf{maximize the total
value} carried out of the knapsack.}

\begin{patternbox}
The single word \emph{``fraction''} (or \emph{``can be broken into
smaller units''}) attached to a knapsack-style capacity problem is the
loudest signal in this entire chapter: it rules out dynamic programming
immediately. The $0/1$ knapsack needs DP precisely because an item is
all-or-nothing, which creates genuine overlapping subproblems (``did we take
item $i$ or not''). The moment items become divisible, that branching
disappears -- we can always characterize the best item to put in next using a
single number, its \emph{value-per-unit-weight ratio}, and greedily fill the
knapsack from the best ratio downward. Any time you see ``maximize total
value/profit under a single weight/capacity/budget constraint, items may be
split,'' reach for the ratio-greedy pattern.
\end{patternbox}

\subsection*{Understanding the problem carefully}

Define, for each item $i$, its \emph{rate of return}
\[
  r_i = \frac{\textit{val}[i]}{\textit{wt}[i]}.
\]
This number tells us exactly how much value we earn per unit of weight we
spend if we take item $i$. Because we are allowed to take fractions, taking
item $i$ is like having access to an unlimited (well, limited by
$\textit{wt}[i]$) supply of ``value at rate $r_i$.'' If our only constraint is
a total weight budget $W$, then to maximize value, common sense says: spend
the budget first on the item with the best rate of return, then the next
best, and so on, until the budget runs out -- possibly using only part of the
last item we touch.

\begin{ideabox}[Greedy Choice]
Sort the items in decreasing order of $r_i = \textit{val}[i] /
\textit{wt}[i]$. Process them in this order; for each item, take as much of
it as the remaining capacity allows -- the whole item if it fits, otherwise
exactly the fraction of it that exactly fills the remaining capacity, after
which the knapsack is full and we stop.
\end{ideabox}

\subsection*{Why the greedy choice is correct}

Suppose, for contradiction, that some optimal solution does not follow this
rule -- say it takes a positive amount $a > 0$ of a lower-ratio item $i$
while leaving some amount $b > 0$ of a higher-ratio item $i'$ (with $r_{i'} >
r_i$) untaken, even though there was room for it. Then we can perform a small
exchange: reduce the amount of item $i$ taken by $\varepsilon =
\min(a, b)$ in weight, and increase the amount of item $i'$ taken by the same
weight $\varepsilon$. The total weight used is unchanged, so the solution
remains feasible. But the value changes by
\[
  \Delta = \varepsilon \cdot r_{i'} - \varepsilon \cdot r_i
         = \varepsilon (r_{i'} - r_i) > 0,
\]
since $r_{i'} > r_i$. So the exchange strictly increases total value, which
contradicts the assumption that the original solution was optimal. Hence no
optimal solution can avoid fully preferring higher-ratio items over
lower-ratio ones whenever weight can be moved between them -- which is
exactly the greedy rule. Because fractional amounts are allowed, this
exchange can always be performed down to an infinitesimal degree, so the
greedy strategy is not just \emph{a} good heuristic here, it is provably
optimal, with no gap to close via DP.

\subsection*{Algorithm}

\begin{enumerate}
  \item Compute the ratio $r_i = \textit{val}[i] / \textit{wt}[i]$ for every
        item.
  \item Sort the items in decreasing order of $r_i$.
  \item Initialize $\textit{remaining} \gets W$ and
        $\textit{totalValue} \gets 0$.
  \item For each item $i$ in sorted order:
  \begin{enumerate}
    \item If $\textit{wt}[i] \le \textit{remaining}$: take the whole item.
          $\textit{totalValue} \gets \textit{totalValue} + \textit{val}[i]$;
          $\textit{remaining} \gets \textit{remaining} - \textit{wt}[i]$.
    \item Otherwise: take the fraction
          $f = \textit{remaining} / \textit{wt}[i]$ of this item.
          $\textit{totalValue} \gets \textit{totalValue} + f \cdot
          \textit{val}[i]$; set $\textit{remaining} \gets 0$ and stop --
          the knapsack is now full.
  \end{enumerate}
  \item Return $\textit{totalValue}$.
\end{enumerate}

\begin{algorithm}[H]
\caption{Fractional Knapsack}
\begin{algorithmic}[1]
\Procedure{FractionalKnapsack}{$\textit{val}, \textit{wt}, W$}
  \State sort items by $\textit{val}[i]/\textit{wt}[i]$ in decreasing order
  \State $\textit{remaining} \gets W$, $\textit{totalValue} \gets 0$
  \For{each item $i$ in sorted order}
    \If{$\textit{remaining} = 0$}
      \State \textbf{break}
    \EndIf
    \If{$\textit{wt}[i] \le \textit{remaining}$}
      \State $\textit{totalValue} \gets \textit{totalValue} + \textit{val}[i]$
      \State $\textit{remaining} \gets \textit{remaining} - \textit{wt}[i]$
    \Else
      \State $f \gets \textit{remaining} / \textit{wt}[i]$
      \State $\textit{totalValue} \gets \textit{totalValue} + f \times \textit{val}[i]$
      \State $\textit{remaining} \gets 0$
    \EndIf
  \EndFor
  \State \Return \textit{totalValue}
\EndProcedure
\end{algorithmic}
\end{algorithm}

\subsection*{Dry run}

\dryrun{Take three items with $(\textit{val}, \textit{wt})$ pairs
$(60, 10)$, $(100, 20)$, $(120, 30)$, and knapsack capacity $W = 50$.}

First compute ratios:
\[
  r_1 = 60/10 = 6.0, \qquad r_2 = 100/20 = 5.0, \qquad r_3 = 120/30 = 4.0.
\]
Sorted in decreasing order of ratio, the items are already in the order
$1, 2, 3$.

\begin{itemize}
  \item \textbf{Item 1} ($\textit{val}=60, \textit{wt}=10$): remaining
        capacity is $50$, and $10 \le 50$, so we take all of it.
        $\textit{totalValue} = 60$, $\textit{remaining} = 40$.
  \item \textbf{Item 2} ($\textit{val}=100, \textit{wt}=20$): remaining
        capacity is $40$, and $20 \le 40$, so we take all of it.
        $\textit{totalValue} = 160$, $\textit{remaining} = 20$.
  \item \textbf{Item 3} ($\textit{val}=120, \textit{wt}=30$): remaining
        capacity is $20$, but $30 > 20$, so we cannot take the whole item.
        We take the fraction $f = 20/30 = 2/3$. This contributes
        $\frac{2}{3} \times 120 = 80$ value.
        $\textit{totalValue} = 160 + 80 = 240$, $\textit{remaining} = 0$.
\end{itemize}

The knapsack is now full, so we stop. The maximum value obtainable is
$\boxed{240}$.

\subsection*{C++ implementation}

\begin{lstlisting}
#include <bits/stdc++.h>
using namespace std;

struct Item {
    int value;
    int weight;
};

double fractionalKnapsack(int W, vector<Item>& items) {
    // Sort items by value/weight ratio in decreasing order
    sort(items.begin(), items.end(), [](const Item& a, const Item& b) {
        double r1 = (double)a.value / (double)a.weight;
        double r2 = (double)b.value / (double)b.weight;
        return r1 > r2;
    });

    int remaining = W;
    double totalValue = 0.0;

    for (auto& item : items) {
        if (remaining == 0) break;

        if (item.weight <= remaining) {
            totalValue += item.value;
            remaining -= item.weight;
        } else {
            double fraction = (double)remaining / (double)item.weight;
            totalValue += fraction * item.value;
            remaining = 0;
        }
    }
    return totalValue;
}

int main() {
    vector<Item> items = {{60, 10}, {100, 20}, {120, 30}};
    int W = 50;

    cout << "Maximum value: " << fractionalKnapsack(W, items) << endl;
    return 0;
}
\end{lstlisting}

\begin{complexitybox}
\textbf{Time Complexity.} Sorting the $n$ items by ratio costs
$O(n \log n)$. The subsequent single pass over the sorted items costs
$O(n)$. Hence the total time complexity is
\[
  O(n \log n).
\]
\textbf{Space Complexity.} Sorting is performed in place (or with a small
$O(n)$ auxiliary array depending on implementation), and we use only a
constant number of extra scalar variables, giving auxiliary space
complexity $O(1)$ beyond the input storage (or $O(n)$ if we count the
input array itself, and $O(\log n)$ for the sort's recursion stack).
\end{complexitybox}

\begin{takeawaybox}
Divisibility is the tell. The instant a knapsack-style problem allows
\emph{fractions} of an item, the all-or-nothing branching that forces
dynamic programming in $0/1$ knapsack disappears, and a single
value-per-weight ratio becomes a perfect, provably optimal greedy ranking.
Contrast this sharply with Section~\ref{sec:assign-cookies}: there, sorting
gave us a feasibility-matching greedy; here, sorting by ratio gives us a
value-maximizing greedy, but the underlying exchange-argument justification
is the same family of reasoning.
\end{takeawaybox}

\section{Greedy Algorithm to Find Minimum Number of Coins}
\label{sec:minimum-coins}

\problemstatement{We are given a target amount $V$ and an unlimited supply of
coins (and notes) of standard denominations -- in the version used on the
sheet, the Indian currency denominations
\[
  \{1,\,2,\,5,\,10,\,20,\,50,\,100,\,500,\,1000\}.
\]
We must choose a multiset of these denominations whose values sum exactly to
$V$, using as \textbf{few coins/notes as possible}, and report that minimum
count (or the coins themselves).}

\begin{patternbox}
\emph{``Minimum number of coins/notes to make a given amount,''} when the
denomination set is a small, fixed, real-world currency system, is the
textbook greedy ``always use the largest denomination that fits'' problem.
The crucial, easy-to-miss keyword here is which \emph{specific} denomination
set is being used. If the problem instead says \emph{``given an arbitrary
list of coin denominations''} (as in the famous LeetCode ``Coin Change''
problem), the greedy approach can silently give a wrong answer, and you must
switch to dynamic programming. Recognizing \emph{which} of these two
near-identical-looking problems you are facing is itself the key skill being
tested here.
\end{patternbox}

\subsection*{Understanding the problem carefully}

For a fixed real-world currency system such as the Indian denominations
above, there is a special structural property: each denomination is large
enough, relative to the ones below it, that repeatedly taking the largest
denomination not exceeding the remaining amount never forces us into a
worse total count later. Concretely, every denomination in this list is at
least double the previous one except for the jumps $1 \to 2$, $2 \to 5$,
$5 \to 10$, $10 \to 20$, $20 \to 50$ -- and a short case check (which we will
not belabor here, since it is a property of \emph{this particular} currency
system rather than something we derive from first principles) confirms that
this set is a so-called \emph{canonical coin system}: the greedy algorithm
always produces an optimal (minimum-count) representation for it.

\begin{ideabox}[Greedy Choice]
Sort the denominations in decreasing order. Repeatedly subtract the largest
denomination that does not exceed the remaining amount, counting how many
times each denomination is used, until the remaining amount reaches zero.
\end{ideabox}

\subsection*{Why greedy works here, and why it can fail elsewhere}

The intuitive argument for why greedy is safe here is this: using one coin
of the largest available denomination $d$ removes as much value as possible
in a single coin. If using a smaller denomination instead could ever lead to
a strictly smaller total coin count, it would have to be because the smaller
coins ``line up'' more efficiently with the rest of the amount -- but because
each denomination in this system divides evenly into, or is closely
compatible with, the next one up, no such inefficiency can arise. This is
\emph{not} a universal truth about coin systems; it is a property that must
be verified for the specific set of denominations given.

To see the failure mode, suppose instead the denominations were
$\{1, 3, 4\}$ and the target amount were $V = 6$. Greedy would take one coin
of $4$, leaving $2$, then two coins of $1$, for a total of three coins:
$4+1+1=6$. But the optimal answer uses two coins of $3$: $3+3=6$, a total of
only two coins. Here greedy fails because $4$ is not a ``safe'' large
denomination relative to $3$ -- taking it first blocks the more efficient
combination. This is exactly why the general ``Coin Change: minimum coins''
problem (with an arbitrary denomination set) requires dynamic programming,
while this specific problem, restricted to real-world currency
denominations, is safe for greedy. Always check which version of the
problem you are being asked before reaching for greedy.

\subsection*{Algorithm}

\begin{enumerate}
  \item Store the denominations in decreasing order:
        $1000, 500, 100, 50, 20, 10, 5, 2, 1$.
  \item Initialize $\textit{remaining} \gets V$ and an empty result list.
  \item For each denomination $d$ in decreasing order:
  \begin{enumerate}
    \item While $d \le \textit{remaining}$: append $d$ to the result list,
          and set $\textit{remaining} \gets \textit{remaining} - d$.
  \end{enumerate}
  \item Return the result list (its length is the minimum coin count).
\end{enumerate}

\begin{algorithm}[H]
\caption{Minimum Number of Coins (Canonical System)}
\begin{algorithmic}[1]
\Procedure{MinimumCoins}{$V$}
  \State $\textit{denoms} \gets [1000, 500, 100, 50, 20, 10, 5, 2, 1]$
  \State $\textit{remaining} \gets V$
  \State $\textit{result} \gets$ empty list
  \For{each $d$ in \textit{denoms}}
    \While{$d \le \textit{remaining}$}
      \State append $d$ to \textit{result}
      \State $\textit{remaining} \gets \textit{remaining} - d$
    \EndWhile
  \EndFor
  \State \Return \textit{result}
\EndProcedure
\end{algorithmic}
\end{algorithm}

\subsection*{Dry run}

\dryrun{Take $V = 93$.}

\begin{itemize}
  \item $d = 1000$: $1000 > 93$, skip.
  \item $d = 500$: $500 > 93$, skip.
  \item $d = 100$: $100 > 93$, skip.
  \item $d = 50$: $50 \le 93$. Take one $50$. $\textit{remaining} = 43$.
        $50 \le 43$? No, stop this denomination.
  \item $d = 20$: $20 \le 43$. Take one $20$. $\textit{remaining} = 23$.
        $20 \le 23$. Take another $20$. $\textit{remaining} = 3$.
        $20 \le 3$? No, stop.
  \item $d = 10$: $10 > 3$, skip.
  \item $d = 5$: $5 > 3$, skip.
  \item $d = 2$: $2 \le 3$. Take one $2$. $\textit{remaining} = 1$.
        $2 \le 1$? No, stop.
  \item $d = 1$: $1 \le 1$. Take one $1$. $\textit{remaining} = 0$.
\end{itemize}

The coins used are $\{50, 20, 20, 2, 1\}$, a total of $5$ coins, and indeed
$50 + 20 + 20 + 2 + 1 = 93$. No combination of these denominations can make
$93$ with fewer than $5$ coins.

\subsection*{C++ implementation}

\begin{lstlisting}
#include <bits/stdc++.h>
using namespace std;

vector<int> minimumCoins(int V) {
    vector<int> denoms = {1000, 500, 100, 50, 20, 10, 5, 2, 1};
    vector<int> result;

    int remaining = V;
    for (int d : denoms) {
        while (d <= remaining) {
            result.push_back(d);
            remaining -= d;
        }
    }
    return result;
}

int main() {
    int V = 93;
    vector<int> coins = minimumCoins(V);

    cout << "Number of coins used: " << coins.size() << endl;
    cout << "Coins: ";
    for (int c : coins) cout << c << " ";
    cout << endl;
    return 0;
}
\end{lstlisting}

\begin{complexitybox}
\textbf{Time Complexity.} There are only $9$ fixed denominations, so the
outer loop is $O(1)$ in the number of denominations. The inner
\textbf{while} loop, across all denominations, executes at most once per
coin actually placed into the result, and the number of coins used is
$O(\log V)$ in the worst case for this denomination system (since
denominations grow roughly geometrically). So the overall time complexity is
\[
  O(\log V),
\]
effectively a constant for realistic input sizes.
\textbf{Space Complexity.} We store the result list of coins used, which has
length $O(\log V)$, so the space complexity is $O(\log V)$ for the output
(or $O(1)$ auxiliary space if we only need the \emph{count} rather than the
list of coins itself).
\end{complexitybox}

\begin{takeawaybox}
Greedy coin selection is only correct for \emph{canonical} denomination
systems -- ones where the structure of the denominations guarantees that
taking the largest coin first never blocks an optimal solution. Always
distinguish this fixed-currency version of the problem from the general
``minimum coins from an arbitrary denomination list'' problem, which has no
such guarantee and requires dynamic programming (typically an unbounded
knapsack-style DP over amounts $0, \dots, V$).
\end{takeawaybox}

\section{Lemonade Change}
\label{sec:lemonade-change}

\problemstatement{Customers stand in a queue to buy lemonade, and each glass
costs exactly \$5. We are given an array \textit{bills}, where
$\textit{bills}[i]$ is the bill the $i$-th customer pays with -- always a
\$5, \$10, or \$20 note. We start with no change at all. For every customer,
\textbf{immediately} after they pay, we must hand back the correct change
using bills we are currently holding (we cannot use a bill the same customer
just gave us as their own change). We must process customers strictly in the
given order and decide whether it is possible to give correct change to
\emph{every} customer; return \texttt{true} if so, and \texttt{false} the
moment it becomes impossible.}

\begin{patternbox}
Two phrases give this one away. First, \emph{``customers in a queue,
processed one at a time, in order''} tells us this is an \emph{online
simulation} -- we are never allowed to look ahead or rearrange the input, so
any algorithm that tries to plan globally (sort, search over future choices)
is solving the wrong problem. Second, \emph{``immediately give change''}
combined with only three possible bill values means that at any moment the
entire relevant state of the world is just two small counters: how many \$5
notes and how many \$10 notes we are holding (\$20 notes are never given as
change, so we never need to count them). Whenever a problem reduces to
``maintain a couple of counters and react instantly to each new input,'' it
is a simulation problem with, at most, one local greedy decision per step --
not a search problem.
\end{patternbox}

\subsection*{Understanding the problem carefully}

Since lemonade costs \$5, the change owed depends only on what the customer
pays with:

\begin{itemize}
  \item paying with \$5 $\Rightarrow$ no change owed;
  \item paying with \$10 $\Rightarrow$ owe \$5;
  \item paying with \$20 $\Rightarrow$ owe \$15.
\end{itemize}

A \$5 note can only ever be used to pay back \$5 of change, never \$15
worth by itself, so it has exactly one job: covering a \$10 customer, or
participating in the \$15 change for a \$20 customer. A \$10 note is even
more restricted: it is completely useless for a \$10 customer (we cannot
give a customer back their own denomination as part of correct change
composition here -- effectively, the only practical use of a \$10 note as
change is as part of the \$15 owed to a \$20 customer, paired with one \$5
note). This asymmetry is the heart of the problem: \textbf{\$5 notes are
strictly more flexible than \$10 notes}, because a \$5 note can help pay
back a \$10 customer \emph{or} a \$20 customer, while a \$10 note can only
ever help pay back a \$20 customer.

\subsection*{The only real decision point}

Customers paying with \$5 require no decision -- we simply keep the note.
Customers paying with \$10 also require no decision -- if we have a \$5 note
we must use it (it is the only way to make \$5 change), and if we do not,
failure is immediate and unavoidable, since no other combination of notes we
might be holding sums to exactly \$5 with the unit notes available.

The only point where two different valid strategies exist is a customer
paying with \$20, who is owed \$15. There are exactly two ways to assemble
\$15 from the notes we hold:

\[
  \$10 + \$5 \qquad \text{or} \qquad \$5 + \$5 + \$5.
\]

\begin{ideabox}[Greedy Choice]
At a \$20 customer, if we hold at least one \$10 note and at least one \$5
note, prefer to give \textbf{\$10 + \$5} over \textbf{\$5 + \$5 + \$5}. Only
fall back to spending three \$5 notes if we do not have a spare \$10 note
available.
\end{ideabox}

\subsection*{Why this greedy choice is safe}

Both options return exactly \$15, so the customer is equally satisfied
either way; the difference is entirely in what we are left holding
afterwards. Giving \$10+\$5 costs us one \$5 note. Giving \$5+\$5+\$5 costs
us three \$5 notes. Since \$5 notes are the more broadly useful resource --
the \emph{only} resource capable of covering a future \$10 customer, and a
necessary ingredient (alongside a \$10 note) for covering a future \$20
customer -- depleting three of them at once when we could have depleted only
one is never better for the future, and is sometimes strictly worse: it
might leave us without enough \$5 notes to cover a \$10 customer who arrives
later. Formally, suppose an alternative strategy chooses \$5+\$5+\$5 for some
\$20 customer while a \$10 note sits unused in hand. We can always replace
that choice with \$10+\$5 instead, ending with two extra \$5 notes in hand
compared to the alternative strategy at every subsequent point in the queue,
and strictly more \$5 notes in hand can never cause a future failure that
would not also have occurred with fewer -- so the greedy choice dominates.
\$10 notes, by contrast, have no other use, so spending one whenever
possible is pure upside.

\subsection*{Algorithm}

\begin{enumerate}
  \item Initialize $\textit{five} \gets 0$, $\textit{ten} \gets 0$.
  \item For each bill $b$ in \textit{bills}, in order:
  \begin{enumerate}
    \item If $b = 5$: $\textit{five} \gets \textit{five} + 1$.
    \item If $b = 10$: if $\textit{five} = 0$, return \texttt{false};
          otherwise $\textit{five} \gets \textit{five} - 1$ and
          $\textit{ten} \gets \textit{ten} + 1$.
    \item If $b = 20$:
    \begin{itemize}
      \item if $\textit{ten} \ge 1$ and $\textit{five} \ge 1$: subtract one
            from each;
      \item else if $\textit{five} \ge 3$: subtract three from
            \textit{five};
      \item else: return \texttt{false}.
    \end{itemize}
  \end{enumerate}
  \item If every customer was processed without returning \texttt{false},
        return \texttt{true}.
\end{enumerate}

\begin{algorithm}[H]
\caption{Lemonade Change}
\begin{algorithmic}[1]
\Procedure{LemonadeChange}{\textit{bills}}
  \State $\textit{five} \gets 0$, $\textit{ten} \gets 0$
  \For{each $b$ in \textit{bills}}
    \If{$b = 5$}
      \State $\textit{five} \gets \textit{five} + 1$
    \ElsIf{$b = 10$}
      \If{$\textit{five} = 0$} \Return \textbf{false} \EndIf
      \State $\textit{five} \gets \textit{five} - 1$
      \State $\textit{ten} \gets \textit{ten} + 1$
    \Else \Comment{$b = 20$}
      \If{$\textit{ten} \ge 1$ \textbf{and} $\textit{five} \ge 1$}
        \State $\textit{ten} \gets \textit{ten} - 1$
        \State $\textit{five} \gets \textit{five} - 1$
      \ElsIf{$\textit{five} \ge 3$}
        \State $\textit{five} \gets \textit{five} - 3$
      \Else
        \State \Return \textbf{false}
      \EndIf
    \EndIf
  \EndFor
  \State \Return \textbf{true}
\EndProcedure
\end{algorithmic}
\end{algorithm}

\subsection*{Dry run}

\dryrun{Take $\textit{bills} = [5, 5, 5, 10, 20]$.}

Start with $\textit{five}=0$, $\textit{ten}=0$.

\begin{itemize}
  \item Bill $5$: $\textit{five}=1$.
  \item Bill $5$: $\textit{five}=2$.
  \item Bill $5$: $\textit{five}=3$.
  \item Bill $10$: $\textit{five} > 0$, so give one \$5. $\textit{five}=2$,
        $\textit{ten}=1$.
  \item Bill $20$: $\textit{ten}\ge1$ and $\textit{five}\ge1$, so give
        \$10+\$5. $\textit{five}=1$, $\textit{ten}=0$.
\end{itemize}

All five customers are served, so the answer is \texttt{true}.

Now take $\textit{bills} = [5, 5, 10, 10, 20]$.

\begin{itemize}
  \item Bill $5$: $\textit{five}=1$.
  \item Bill $5$: $\textit{five}=2$.
  \item Bill $10$: give one \$5. $\textit{five}=1$, $\textit{ten}=1$.
  \item Bill $10$: give one \$5. $\textit{five}=0$, $\textit{ten}=2$.
  \item Bill $20$: need $\textit{ten}\ge1$ and $\textit{five}\ge1$, but
        $\textit{five}=0$. Fall back to $\textit{five}\ge3$, but again
        $\textit{five}=0$. Neither option works, so return \texttt{false}.
\end{itemize}

\subsection*{C++ implementation}

\begin{lstlisting}
#include <bits/stdc++.h>
using namespace std;

bool lemonadeChange(vector<int>& bills) {
    int five = 0, ten = 0;

    for (int bill : bills) {
        if (bill == 5) {
            five++;
        } else if (bill == 10) {
            if (five == 0) return false;
            five--;
            ten++;
        } else { // bill == 20
            if (ten >= 1 && five >= 1) {
                ten--;
                five--;
            } else if (five >= 3) {
                five -= 3;
            } else {
                return false;
            }
        }
    }
    return true;
}

int main() {
    vector<int> bills = {5, 5, 5, 10, 20};
    cout << (lemonadeChange(bills) ? "true" : "false") << endl;
    return 0;
}
\end{lstlisting}

\begin{complexitybox}
\textbf{Time Complexity.} Each of the $n$ customers is processed exactly
once, with only constant work (a couple of comparisons and increments) per
customer, giving
\[
  O(n).
\]
\textbf{Space Complexity.} We track only the two counters \textit{five} and
\textit{ten}, so the space complexity is $O(1)$.
\end{complexitybox}

\begin{takeawaybox}
When several local choices return the same immediate result (here, both
change-giving options return exactly \$15), break the tie by asking which
choice preserves the resource that is more broadly useful for the future.
Here, \$5 notes serve two future needs while \$10 notes serve only one, so
spending a \$10 note whenever possible is always at least as good, and
sometimes strictly necessary, for surviving the rest of the queue.
\end{takeawaybox}

\section{Valid Parenthesis String}
\label{sec:valid-parenthesis-string}

\problemstatement{We are given a string $s$ consisting only of the
characters \texttt{'('}, \texttt{')'}, and \texttt{'*'}. The character
\texttt{'*'} is a wildcard: each occurrence of it may be treated, independently,
as either \texttt{'('}, or \texttt{')'}, or the empty string. We must decide
whether there is \emph{some} way to resolve every \texttt{'*'} so that the
resulting string is a valid parenthesis sequence (every prefix has at least
as many \texttt{'('} as \texttt{')'}, and the total counts of \texttt{'('}
and \texttt{')'} are equal).}

\begin{patternbox}
The phrase \emph{``each wildcard may independently be one of several
things; does \textbf{some} assignment make the whole thing valid''} looks,
at first glance, like it demands trying all $3^k$ assignments of the $k$
wildcards -- an exponential brute force. But the underlying validity
condition for parentheses is itself expressible as a simple running balance
(``open count never goes negative, and ends at zero''), and whenever a
validity condition can be tracked by a \emph{running numeric range} rather
than an exact value, that is a strong signal that we can collapse all
exponentially many choices into tracking just the minimum and maximum
possible value of that running quantity at each step -- a classic
``track an interval of reachable states'' greedy, no recursion or
backtracking required.
\end{patternbox}

\subsection*{Understanding the problem carefully}

For a string with only \texttt{'('} and \texttt{')'}, validity is checked by
maintaining a single counter, the \emph{open balance}: add $1$ for
\texttt{'('}, subtract $1$ for \texttt{')'}, and the string is valid exactly
when the balance never goes negative and ends at exactly $0$.

With \texttt{'*'} in the mix, the open balance after processing a prefix is
no longer a single number -- it is a \emph{range} of possible numbers,
depending on how each \texttt{'*'} so far was resolved. Let
$\textit{lo}$ be the smallest possible balance consistent with some
resolution of the wildcards seen so far, and $\textit{hi}$ be the largest
possible balance. We update this range character by character:

\begin{itemize}
  \item On \texttt{'('}: both the smallest and largest possible balance
        increase by $1$, so $\textit{lo} \mathrel{+}= 1$,
        $\textit{hi} \mathrel{+}= 1$.
  \item On \texttt{')'}: both decrease by $1$, so
        $\textit{lo} \mathrel{-}= 1$, $\textit{hi} \mathrel{-}= 1$.
  \item On \texttt{'*'}: in the most pessimistic resolution (treat it as
        \texttt{')'}), the balance decreases by $1$, so
        $\textit{lo} \mathrel{-}= 1$. In the most optimistic resolution
        (treat it as \texttt{'('}), the balance increases by $1$, so
        $\textit{hi} \mathrel{+}= 1$. (Treating it as empty would leave the
        balance unchanged, which always lies between these two extremes, so
        it never needs separate tracking.)
\end{itemize}

\begin{ideabox}[Greedy Choice]
At every position, maintain only the \emph{interval} $[\textit{lo},
\textit{hi}]$ of balances that are achievable by \emph{some} resolution of
the wildcards seen so far -- not every individual achievable value, just the
two extremes. Whenever $\textit{hi}$ drops below $0$, no resolution can
recover (the balance has gone irrecoverably negative even in the best case),
so we can answer \texttt{false} immediately. Whenever $\textit{lo}$ drops
below $0$, clamp it to $0$, because a balance can never validly be negative,
so any resolution that would have produced a negative running balance is
simply discarded from consideration -- it is never the right choice for an
eventual valid string, so we do not need to track it.
\end{ideabox}

\subsection*{Why tracking only the interval is enough}

It might seem like we are throwing away information by tracking only
$[\textit{lo}, \textit{hi}]$ rather than the full set of achievable balances.
But notice that the set of achievable balances after processing a prefix is
always a contiguous range of integers with the \emph{same parity step of
$1$} apart (because every character changes the balance by exactly $-1$,
$0$, or $+1$, and a \texttt{'*'} contributes all three options, which fills
in every integer between the all-\texttt{')'} extreme and the all-\texttt{'('}
extreme). So the reachable set truly is an interval, and tracking its
endpoints loses nothing. At the very end, the string can be made valid if
and only if $0$ lies in the final interval -- but since we always clamp
$\textit{lo}$ to be at least $0$ throughout, checking that the final
$\textit{lo} = 0$ is equivalent to checking that $0$ is reachable, because if
$0$ were reachable, clamping would have kept $\textit{lo}$ pinned at (or
below, then clamped to) $0$ rather than drifting strictly positive.

\subsection*{Algorithm}

\begin{enumerate}
  \item Initialize $\textit{lo} \gets 0$, $\textit{hi} \gets 0$.
  \item For each character $c$ in $s$:
  \begin{enumerate}
    \item If $c = $ \texttt{'('}: $\textit{lo} \gets \textit{lo}+1$,
          $\textit{hi} \gets \textit{hi}+1$.
    \item If $c = $ \texttt{')'}: $\textit{lo} \gets \textit{lo}-1$,
          $\textit{hi} \gets \textit{hi}-1$.
    \item If $c = $ \texttt{'*'}: $\textit{lo} \gets \textit{lo}-1$,
          $\textit{hi} \gets \textit{hi}+1$.
    \item If $\textit{hi} < 0$: return \texttt{false} immediately.
    \item If $\textit{lo} < 0$: clamp $\textit{lo} \gets 0$.
  \end{enumerate}
  \item After processing the whole string, return
        $(\textit{lo} = 0)$.
\end{enumerate}

\begin{algorithm}[H]
\caption{Valid Parenthesis String}
\begin{algorithmic}[1]
\Procedure{CheckValidString}{$s$}
  \State $\textit{lo} \gets 0$, $\textit{hi} \gets 0$
  \For{each character $c$ in $s$}
    \If{$c = $ '('}
      \State $\textit{lo} \gets \textit{lo} + 1$, $\textit{hi} \gets \textit{hi} + 1$
    \ElsIf{$c = $ ')'}
      \State $\textit{lo} \gets \textit{lo} - 1$, $\textit{hi} \gets \textit{hi} - 1$
    \Else \Comment{$c = $ '*'}
      \State $\textit{lo} \gets \textit{lo} - 1$, $\textit{hi} \gets \textit{hi} + 1$
    \EndIf
    \If{$\textit{hi} < 0$}
      \State \Return \textbf{false}
    \EndIf
    \If{$\textit{lo} < 0$}
      \State $\textit{lo} \gets 0$
    \EndIf
  \EndFor
  \State \Return $(\textit{lo} = 0)$
\EndProcedure
\end{algorithmic}
\end{algorithm}

\subsection*{Dry run}

\dryrun{Take $s = \texttt{"(*))"}$.}

Start with $\textit{lo}=0$, $\textit{hi}=0$.

\begin{itemize}
  \item $c = $ \texttt{'('}: $\textit{lo}=1$, $\textit{hi}=1$.
  \item $c = $ \texttt{'*'}: $\textit{lo}=0$, $\textit{hi}=2$.
  \item $c = $ \texttt{')'}: $\textit{lo}=-1 \to$ clamp to $0$,
        $\textit{hi}=1$.
  \item $c = $ \texttt{')'}: $\textit{lo}=-1 \to$ clamp to $0$,
        $\textit{hi}=0$.
\end{itemize}

At no point did $\textit{hi}$ go negative, and the final $\textit{lo} = 0$,
so the answer is \texttt{true}. We can confirm this directly: resolving the
\texttt{'*'} as \texttt{'('} turns $s$ into \texttt{"(())"}, whose balances
are $1,2,1,0$ -- never negative, ending at $0$ -- so this is indeed a valid
resolution, and the algorithm's \texttt{true} verdict is justified.

Now take $s = \texttt{"(*))("}$, which appends one extra unmatched
\texttt{'('} to the previous string.

\begin{itemize}
  \item As above, after \texttt{"(*))"}: $\textit{lo}=0$, $\textit{hi}=0$.
  \item $c = $ \texttt{'('}: $\textit{lo}=1$, $\textit{hi}=1$.
\end{itemize}

Final $\textit{lo} = 1 \neq 0$, so the answer is \texttt{false} -- correctly,
since the trailing unmatched \texttt{'('} can never be closed by anything
after it.

\subsection*{C++ implementation}

\begin{lstlisting}
#include <bits/stdc++.h>
using namespace std;

bool checkValidString(string s) {
    int lo = 0, hi = 0;

    for (char c : s) {
        if (c == '(') {
            lo++;
            hi++;
        } else if (c == ')') {
            lo--;
            hi--;
        } else { // c == '*'
            lo--;
            hi++;
        }

        if (hi < 0) return false;
        if (lo < 0) lo = 0;
    }
    return lo == 0;
}

int main() {
    string s = "(*))";
    cout << (checkValidString(s) ? "true" : "false") << endl;
    return 0;
}
\end{lstlisting}

\begin{complexitybox}
\textbf{Time Complexity.} We scan the string once, doing constant work per
character, giving
\[
  O(n).
\]
\textbf{Space Complexity.} We maintain only the two integers \textit{lo} and
\textit{hi}, so the space complexity is $O(1)$.
\end{complexitybox}

\begin{takeawaybox}
Whenever a validity check can be expressed as a single running numeric
quantity that must stay within bounds, and the input introduces wildcards
with several possible numeric effects, try tracking the \emph{interval} of
reachable values for that quantity instead of branching over every
combination of wildcard resolutions. This turns an exponential
``try every assignment'' search into a single linear greedy sweep.
\end{takeawaybox}

\section{N Meetings in One Room}
\label{sec:n-meetings}

\problemstatement{We are given $n$ meetings, where the $i$-th meeting has a
start time $\textit{start}[i]$ and a finish time $\textit{end}[i]$. Only one
meeting room is available, so a meeting occupies the room for its entire
duration $[\textit{start}[i], \textit{end}[i]]$, and the room becomes free
only strictly after $\textit{end}[i]$ -- so two meetings can both be held in
the room only if one of them starts strictly after the other has ended. We
must select the \textbf{maximum number of non-overlapping meetings} that can
be held in this single room.}

\begin{patternbox}
\emph{``One resource (a room), many requests each with a start and an end,
maximize the number of non-conflicting requests you can serve''} is the
exact signature of the classical \emph{Activity Selection Problem}, the
flagship example of interval scheduling. The keyword combination to listen
for is \emph{``maximum number of non-overlapping intervals''} together with
\emph{a single shared resource}. The unintuitive but crucial insight that
makes this greedy rather than a search problem is that you should sort by
\textbf{finish time}, not by start time and not by duration -- a fact that
is worth memorizing because it differs from a tempting but incorrect first
instinct.
\end{patternbox}

\subsection*{Understanding the problem carefully}

It is tempting to think that scheduling the \emph{shortest} meetings first,
or the \emph{earliest-starting} meetings first, would leave the most room for
others. Both of these intuitions are wrong, and small examples expose the
flaw immediately. Consider two meetings, $A = [1, 100]$ and $B = [2, 3]$.
Sorting by start time would consider $A$ first (it starts earliest); taking
$A$ blocks $B$ entirely, leaving only $1$ meeting scheduled, whereas taking
$B$ instead would have left room afterward for further meetings starting
after time $3$. Sorting by duration would favor $B$ here by luck, but that
is not a reliable rule either -- a short meeting that starts very early and
a slightly longer meeting that finishes only one unit later can still flip
which choice is better, depending on what comes after.

What actually matters is not how long a meeting takes, nor when it starts,
but \emph{when it frees up the room for everyone else}. A meeting that
finishes early frees the room early, leaving the largest possible remaining
time window for future meetings -- regardless of how late it started or how
long it lasted.

\begin{ideabox}[Greedy Choice]
Sort all meetings by finish time, in increasing order. Always select the
next meeting (in this sorted order) whose start time is strictly after the
finish time of the most recently selected meeting.
\end{ideabox}

\subsection*{Why sorting by finish time is correct}

This is the standard exchange-argument proof used for activity selection.
Let $m_1$ be the meeting with the earliest finish time among all $n$
meetings. We claim there is always an optimal solution that includes $m_1$.
Suppose some optimal solution $O$ does not include $m_1$, but includes some
other meeting $m'$ as its first (earliest-finishing) chosen meeting instead.
Since $m_1$ has the globally earliest finish time of \emph{all} meetings, in
particular $\textit{end}[m_1] \le \textit{end}[m']$. Replacing $m'$ with
$m_1$ inside $O$ cannot create any new conflict with the meetings that come
after $m'$ in $O$, because $m_1$ frees the room \emph{no later} than $m'$
did. So this swap produces another solution that is still feasible and has
exactly the same number of meetings as $O$ -- meaning some optimal solution
does include $m_1$. Once $m_1$ is fixed as selected, the remaining problem
is exactly the same problem, restricted to meetings that start after
$\textit{end}[m_1]$, so the same argument applies recursively (or, in
iterative form, repeatedly) to the next-earliest-finishing remaining
meeting, and so on. This is what justifies scanning the finish-time-sorted
list exactly once, greedily picking whichever next meeting is compatible
with the last one picked.

\subsection*{Algorithm}

\begin{enumerate}
  \item Sort all $n$ meetings by $\textit{end}[i]$ in increasing order.
  \item Initialize $\textit{count} \gets 1$ and
        $\textit{lastEnd} \gets \textit{end}$ of the first meeting in
        sorted order (the very first meeting is always selected, since
        nothing has been scheduled yet to conflict with it).
  \item For each remaining meeting $i$ in sorted order, starting from the
        second one:
  \begin{enumerate}
    \item If $\textit{start}[i] > \textit{lastEnd}$: select this meeting.
          $\textit{count} \gets \textit{count} + 1$;
          $\textit{lastEnd} \gets \textit{end}[i]$.
    \item Otherwise: skip this meeting (it conflicts with the last
          selected meeting).
  \end{enumerate}
  \item Return \textit{count}.
\end{enumerate}

\begin{algorithm}[H]
\caption{Maximum Meetings in One Room}
\begin{algorithmic}[1]
\Procedure{MaxMeetings}{$\textit{start}, \textit{end}$}
  \State sort meetings by $\textit{end}[i]$ in increasing order
  \State $\textit{count} \gets 1$
  \State $\textit{lastEnd} \gets \textit{end}[0]$ \Comment{always take the first meeting}
  \For{$i \gets 1$ to $n-1$}
    \If{$\textit{start}[i] > \textit{lastEnd}$}
      \State $\textit{count} \gets \textit{count} + 1$
      \State $\textit{lastEnd} \gets \textit{end}[i]$
    \EndIf
  \EndFor
  \State \Return \textit{count}
\EndProcedure
\end{algorithmic}
\end{algorithm}

\subsection*{Dry run}

\dryrun{Take meetings with $(\textit{start}, \textit{end})$ pairs:
$(1,2), (3,4), (0,6), (5,7), (8,9), (5,9)$.}

Sorting by finish time gives the order:
\[
  (1,2),\ (3,4),\ (0,6),\ (5,7),\ (8,9),\ (5,9).
\]

\begin{itemize}
  \item Select $(1,2)$ unconditionally. $\textit{count}=1$,
        $\textit{lastEnd}=2$.
  \item $(3,4)$: $\textit{start}=3 > \textit{lastEnd}=2$. Select.
        $\textit{count}=2$, $\textit{lastEnd}=4$.
  \item $(0,6)$: $\textit{start}=0 \not> \textit{lastEnd}=4$. Skip.
  \item $(5,7)$: $\textit{start}=5 > \textit{lastEnd}=4$. Select.
        $\textit{count}=3$, $\textit{lastEnd}=7$.
  \item $(8,9)$: $\textit{start}=8 > \textit{lastEnd}=7$. Select.
        $\textit{count}=4$, $\textit{lastEnd}=9$.
  \item $(5,9)$: $\textit{start}=5 \not> \textit{lastEnd}=9$. Skip.
\end{itemize}

The maximum number of meetings that can be held is $\boxed{4}$, namely
$(1,2), (3,4), (5,7), (8,9)$.

\subsection*{C++ implementation}

\begin{lstlisting}
#include <bits/stdc++.h>
using namespace std;

int maxMeetings(vector<int>& start, vector<int>& end) {
    int n = start.size();
    vector<pair<int,int>> meetings(n);
    for (int i = 0; i < n; i++) {
        meetings[i] = {end[i], start[i]};
    }

    sort(meetings.begin(), meetings.end()); // sorts by end time first

    int count = 1;
    int lastEnd = meetings[0].first;

    for (int i = 1; i < n; i++) {
        int s = meetings[i].second;
        int e = meetings[i].first;
        if (s > lastEnd) {
            count++;
            lastEnd = e;
        }
    }
    return count;
}

int main() {
    vector<int> start = {1, 3, 0, 5, 8, 5};
    vector<int> end   = {2, 4, 6, 7, 9, 9};

    cout << "Maximum meetings: " << maxMeetings(start, end) << endl;
    return 0;
}
\end{lstlisting}

\begin{complexitybox}
\textbf{Time Complexity.} Sorting the $n$ meetings by finish time costs
$O(n \log n)$. The subsequent single linear pass costs $O(n)$. So the
overall time complexity is
\[
  O(n \log n).
\]
\textbf{Space Complexity.} We use an auxiliary array of $n$ pairs to sort
by finish time, giving $O(n)$ space (this can be reduced to $O(1)$
auxiliary space if the input is given as an array of pairs already, sorted
in place).
\end{complexitybox}

\begin{takeawaybox}
For interval-scheduling problems with a single shared resource, always sort
by \textbf{finish time}, never by start time or duration. The meeting that
frees the resource soonest is always at least as good a first choice as any
other meeting, by a direct exchange argument -- and this single fact is the
entire content of the activity selection problem, a pattern that reappears
throughout the rest of this chapter under different names.
\end{takeawaybox}

\section{Jump Game}
\label{sec:jump-game}

\problemstatement{We are given an array $\textit{nums}$ of length $n$, where
$\textit{nums}[i]$ is the maximum number of steps we are allowed to jump
forward from index $i$ (we may jump anywhere from $1$ up to
$\textit{nums}[i]$ steps, or choose to jump fewer). Starting at index $0$, we
must decide whether it is possible to reach index $n-1$ \textbf{at all},
following any sequence of valid jumps. Return \texttt{true} or
\texttt{false}.}

\begin{patternbox}
\emph{``Can you reach the last index, given a maximum jump length at each
position''} is the defining example of a \emph{reachability-horizon}
greedy. The keyword to notice is that we only need a yes/no feasibility
answer, not the actual sequence of jumps or the minimum number of jumps --
that distinction matters, because it means we never need to track \emph{how}
we got somewhere, only the single number \emph{how far forward can we
possibly see from everything visited so far}. Whenever a problem asks
``can position $X$ be reached'' under a rule where each position extends
your reach by some amount, track a running \emph{farthest reachable index}
and sweep once, left to right.
\end{patternbox}

\subsection*{Understanding the problem carefully}

Define $\textit{farthest}$ as the largest index reachable using any
combination of jumps from indices we have already confirmed are reachable,
considering only indices up to the one we are currently examining. As we
scan left to right through index $i = 0, 1, \dots, n-1$:

\begin{itemize}
  \item If $i \le \textit{farthest}$, then index $i$ itself is reachable
        (it lies within the horizon established by earlier indices), so we
        are allowed to jump from it: $\textit{farthest}$ can potentially be
        extended to $i + \textit{nums}[i]$.
  \item If $i > \textit{farthest}$, then index $i$ is \emph{not} reachable
        by anything before it, and since we are scanning strictly left to
        right, nothing later can make index $i$ reachable either -- the
        whole array up to this point has failed to bridge the gap, so the
        answer is immediately \texttt{false}.
\end{itemize}

\begin{ideabox}[Greedy Choice]
Maintain a single running value $\textit{farthest}$, the best (largest)
index reachable using everything examined so far. At each index $i$, first
check whether $i$ itself is still within $\textit{farthest}$ (if not, fail
immediately); otherwise, update
$\textit{farthest} \gets \max(\textit{farthest},\, i + \textit{nums}[i])$.
If at any point $\textit{farthest} \ge n-1$, the last index is reachable and
we may stop early and answer \texttt{true}.
\end{ideabox}

\subsection*{Why tracking only the farthest reach is correct}

This works because reachability is monotonic in a useful sense: if index
$j \le \textit{farthest}$ is reachable, every index from $0$ up to
$\textit{farthest}$ is also reachable, because we can always choose to jump
\emph{short} of our maximum allowance at any earlier reachable index to land
exactly where we want (jumping any number of steps from $1$ up to the
maximum is allowed, so every intermediate landing spot is achievable, not
just the farthest one). This means we never need to track the \emph{set} of
reachable indices, nor which specific path of jumps was used -- only the
single supremum of that set, $\textit{farthest}$, fully describes everything
we need to know about reachability at each step. There is no benefit to
"saving" a long jump for later or being clever about which specific jump
length to use, because every index up to the maximum reach is already
reachable regardless of which path got us there. This collapses what could
look like an exponential search over jump sequences into one linear pass
that only ever needs to remember one number.

\subsection*{Algorithm}

\begin{enumerate}
  \item Initialize $\textit{farthest} \gets 0$.
  \item For $i \gets 0$ to $n-1$:
  \begin{enumerate}
    \item If $i > \textit{farthest}$: return \texttt{false} (index $i$, and
          therefore the rest of the array, is unreachable).
    \item $\textit{farthest} \gets \max(\textit{farthest},\, i +
          \textit{nums}[i])$.
    \item If $\textit{farthest} \ge n-1$: return \texttt{true} early.
  \end{enumerate}
  \item If the loop completes without returning \texttt{false}, return
        \texttt{true}.
\end{enumerate}

\begin{algorithm}[H]
\caption{Jump Game (Reachability)}
\begin{algorithmic}[1]
\Procedure{CanJump}{\textit{nums}}
  \State $\textit{farthest} \gets 0$
  \State $n \gets \text{length}(\textit{nums})$
  \For{$i \gets 0$ to $n-1$}
    \If{$i > \textit{farthest}$}
      \State \Return \textbf{false}
    \EndIf
    \State $\textit{farthest} \gets \max(\textit{farthest}, i + \textit{nums}[i])$
    \If{$\textit{farthest} \ge n - 1$}
      \State \Return \textbf{true}
    \EndIf
  \EndFor
  \State \Return \textbf{true}
\EndProcedure
\end{algorithmic}
\end{algorithm}

\subsection*{Dry run}

\dryrun{Take $\textit{nums} = [2, 3, 1, 1, 4]$ ($n=5$, last index $=4$).}

\begin{itemize}
  \item $i=0$: $0 \le \textit{farthest}=0$. Update
        $\textit{farthest} = \max(0, 0+2) = 2$.
  \item $i=1$: $1 \le 2$. Update
        $\textit{farthest} = \max(2, 1+3) = 4$. Since $4 \ge n-1 = 4$,
        return \texttt{true} immediately.
\end{itemize}

The answer is \texttt{true}: from index $0$ we can jump $2$ steps to index
$2$... but more directly, from index $1$ (reachable, since $1 \le 2$) a jump
of $3$ steps reaches index $4$, the last index.

Now take $\textit{nums} = [3, 2, 1, 0, 4]$ ($n=5$, last index $=4$).

\begin{itemize}
  \item $i=0$: $0 \le 0$. $\textit{farthest} = \max(0, 0+3) = 3$.
  \item $i=1$: $1 \le 3$. $\textit{farthest} = \max(3, 1+2) = 3$.
  \item $i=2$: $2 \le 3$. $\textit{farthest} = \max(3, 2+1) = 3$.
  \item $i=3$: $3 \le 3$. $\textit{farthest} = \max(3, 3+0) = 3$.
  \item $i=4$: $4 > \textit{farthest}=3$. Return \texttt{false}.
\end{itemize}

The answer is \texttt{false}: index $3$ has a jump length of $0$, and
everything before it can reach at most index $3$, so index $4$ -- and the
useful $4$ sitting right there at index $4$ -- can never be reached.

\subsection*{C++ implementation}

\begin{lstlisting}
#include <bits/stdc++.h>
using namespace std;

bool canJump(vector<int>& nums) {
    int n = nums.size();
    int farthest = 0;

    for (int i = 0; i < n; i++) {
        if (i > farthest) return false;
        farthest = max(farthest, i + nums[i]);
        if (farthest >= n - 1) return true;
    }
    return true;
}

int main() {
    vector<int> nums = {2, 3, 1, 1, 4};
    cout << (canJump(nums) ? "true" : "false") << endl;
    return 0;
}
\end{lstlisting}

\begin{complexitybox}
\textbf{Time Complexity.} We scan the array once, doing constant work per
index, so the time complexity is
\[
  O(n).
\]
\textbf{Space Complexity.} We maintain only the single variable
\textit{farthest}, so the space complexity is $O(1)$.
\end{complexitybox}

\begin{takeawaybox}
When a problem only asks \emph{whether} some destination is reachable under
a rule that extends your reach from any currently-reachable position, you
rarely need to enumerate paths. Track the single furthest horizon reachable
so far and sweep once -- this is strictly more efficient than recursion or
BFS/DP over individual positions, and it generalizes directly to the
minimum-jumps variant in the next section.
\end{takeawaybox}

\section{Jump Game II}
\label{sec:jump-game-2}

\problemstatement{We are given an array $\textit{nums}$ of length $n$, where
$\textit{nums}[i]$ is the maximum number of steps we may jump forward from
index $i$. It is guaranteed that index $n-1$ is reachable from index $0$.
Starting at index $0$, we must find the \textbf{minimum number of jumps}
required to reach index $n-1$.}

\begin{patternbox}
This is the natural sequel to Section~\ref{sec:jump-game}: the keywords
change from \emph{``is it possible''} to \emph{``minimum number of
jumps,''} which upgrades the problem from a feasibility check to an
optimization. Whenever a reachability-horizon greedy problem is upgraded to
ask for the \emph{minimum number of steps/jumps/levels} to get somewhere,
that is a strong signal to think in terms of \emph{BFS levels} -- here,
each ``level'' is exactly one jump, and we want to know how many levels it
takes for the frontier of reachable positions to engulf the last index. We
will see that this BFS-by-levels idea can be implemented without an actual
queue, using two horizon variables instead.
\end{patternbox}

\subsection*{Understanding the problem carefully}

Think of the jumps as expanding in waves, exactly like breadth-first search
on a graph where node $i$ has an edge to every node in
$[i+1, i+\textit{nums}[i]]$. After $0$ jumps, we have only reached index $0$.
After $1$ jump, we have reached every index reachable directly from index
$0$ -- a contiguous range $[1, \textit{nums}[0]]$. After $2$ jumps, we have
reached every index reachable from \emph{any} index in that first range, and
so on. The number of jumps needed to first include index $n-1$ in the
reachable range is exactly the answer we want.

The key realization is that we do not need an actual BFS queue holding every
individual index. Instead, at each ``wave'' (each fixed number of jumps used
so far), we only need to know the boundary of the current range, and the
farthest index reachable by spending one more jump from anywhere inside the
current range.

\begin{ideabox}[Greedy Choice]
Maintain two horizons: $\textit{currentEnd}$, the farthest index reachable
using the jumps counted so far, and $\textit{farthest}$, the farthest index
reachable using \emph{one more} jump from anywhere within the current
range. Scan left to right; for every index $i$ within the current range,
update $\textit{farthest} \gets \max(\textit{farthest}, i +
\textit{nums}[i])$. The instant we finish processing the current range (when
$i$ reaches $\textit{currentEnd}$), we have exhausted everything reachable
with the current jump count, so we commit to spending one more jump:
increment the jump counter and set $\textit{currentEnd} \gets
\textit{farthest}$, beginning the next wave.
\end{ideabox}

\subsection*{Why this greedy expansion is correct}

This is exactly BFS, level by level, except that instead of storing every
node at the current level in a queue, we only ever need the single number
$\textit{farthest}$ -- the best possible reach from \emph{any} node at the
current level -- because, just as in Jump Game~I, every index between the
previous horizon and the new farthest reach is automatically included in
this level's frontier the moment any node within the previous range can
reach it. We never need to know \emph{which} specific index within the
current range produced the best $\textit{farthest}$ value, only the value
itself, since the next wave starts from the entire new range regardless of
which earlier index unlocked it. This is what allows an $O(n)$ single pass
to replace what would otherwise look like an explicit multi-source BFS.

\subsection*{Algorithm}

\begin{enumerate}
  \item Initialize $\textit{jumps} \gets 0$, $\textit{currentEnd} \gets 0$,
        $\textit{farthest} \gets 0$.
  \item For $i \gets 0$ to $n-2$ (we never need to jump \emph{from} the last
        index):
  \begin{enumerate}
    \item $\textit{farthest} \gets \max(\textit{farthest},\, i +
          \textit{nums}[i])$.
    \item If $i = \textit{currentEnd}$: this wave is exhausted, so spend one
          more jump. $\textit{jumps} \gets \textit{jumps} + 1$;
          $\textit{currentEnd} \gets \textit{farthest}$.
  \end{enumerate}
  \item Return \textit{jumps}.
\end{enumerate}

\begin{algorithm}[H]
\caption{Jump Game II (Minimum Jumps)}
\begin{algorithmic}[1]
\Procedure{MinJumps}{\textit{nums}}
  \State $n \gets \text{length}(\textit{nums})$
  \State $\textit{jumps} \gets 0$, $\textit{currentEnd} \gets 0$, $\textit{farthest} \gets 0$
  \For{$i \gets 0$ to $n-2$}
    \State $\textit{farthest} \gets \max(\textit{farthest}, i + \textit{nums}[i])$
    \If{$i = \textit{currentEnd}$}
      \State $\textit{jumps} \gets \textit{jumps} + 1$
      \State $\textit{currentEnd} \gets \textit{farthest}$
    \EndIf
  \EndFor
  \State \Return \textit{jumps}
\EndProcedure
\end{algorithmic}
\end{algorithm}

\subsection*{Dry run}

\dryrun{Take $\textit{nums} = [2, 3, 1, 1, 4]$ ($n=5$, indices $0..4$).}

Start with $\textit{jumps}=0$, $\textit{currentEnd}=0$,
$\textit{farthest}=0$. We loop $i$ from $0$ to $3$ (i.e.\ $n-2=3$).

\begin{itemize}
  \item $i=0$: $\textit{farthest} = \max(0, 0+2) = 2$. Since
        $i = 0 = \textit{currentEnd}$, spend a jump: $\textit{jumps}=1$,
        $\textit{currentEnd} = 2$.
  \item $i=1$: $\textit{farthest} = \max(2, 1+3) = 4$. Since
        $i = 1 \ne \textit{currentEnd}=2$, no new jump yet.
  \item $i=2$: $\textit{farthest} = \max(4, 2+1) = 4$. Since
        $i = 2 = \textit{currentEnd}$, spend a jump: $\textit{jumps}=2$,
        $\textit{currentEnd} = 4$.
  \item $i=3$: $\textit{farthest} = \max(4, 3+1) = 4$. Since
        $i = 3 \ne \textit{currentEnd}=4$, no new jump.
\end{itemize}

The loop ends (we stop before $i=4$, the last index). The answer is
$\textit{jumps} = 2$: jump from index $0$ to index $1$ (using $\textit{nums}[0]=2$
to reach as far as index $1$, the start of the best next wave), then jump
from index $1$ directly to index $4$ using $\textit{nums}[1]=3$.

\subsection*{C++ implementation}

\begin{lstlisting}
#include <bits/stdc++.h>
using namespace std;

int minJumps(vector<int>& nums) {
    int n = nums.size();
    int jumps = 0, currentEnd = 0, farthest = 0;

    for (int i = 0; i < n - 1; i++) {
        farthest = max(farthest, i + nums[i]);
        if (i == currentEnd) {
            jumps++;
            currentEnd = farthest;
        }
    }
    return jumps;
}

int main() {
    vector<int> nums = {2, 3, 1, 1, 4};
    cout << "Minimum jumps: " << minJumps(nums) << endl;
    return 0;
}
\end{lstlisting}

\begin{complexitybox}
\textbf{Time Complexity.} The single loop visits each of the first $n-1$
indices exactly once, doing constant work per index, giving
\[
  O(n).
\]
\textbf{Space Complexity.} We maintain only three scalar variables
($\textit{jumps}$, $\textit{currentEnd}$, $\textit{farthest}$), so the space
complexity is $O(1)$.
\end{complexitybox}

\begin{takeawaybox}
A feasibility-horizon greedy (Section~\ref{sec:jump-game}) upgrades cleanly
into a minimum-steps greedy by reframing it as implicit, level-by-level BFS:
track the boundary of the current wave (\textit{currentEnd}) and the best
reach achievable from anywhere inside that wave (\textit{farthest}); each
time the scan catches up to the current boundary, a new wave -- one more
jump -- begins. This pattern of ``BFS without an explicit queue, using two
horizon pointers'' generalizes to any problem phrased as ``minimum number of
steps to extend reach to a target.''
\end{takeawaybox}

\section{Minimum Number of Platforms Required for a Railway}
\label{sec:minimum-platforms}

\problemstatement{We are given the arrival times $\textit{arr}[0..n-1]$ and
departure times $\textit{dep}[0..n-1]$ of $n$ trains scheduled at a single
railway station, where train $i$ occupies a platform from
$\textit{arr}[i]$ until $\textit{dep}[i]$ (inclusive; if one train departs
at the exact same time another arrives, they are considered to overlap and
require separate platforms). We must find the \textbf{minimum number of
platforms} the station needs so that no train is ever forced to wait for a
free platform.}

\begin{patternbox}
\emph{``Minimum number of resources (platforms/rooms/servers) needed so that
no request ever waits, given many overlapping time intervals''} is the
classic \emph{maximum overlap / interval point sweep} pattern -- it is
closely related to, but distinct from, the activity-selection pattern in
Section~\ref{sec:n-meetings}. There, we had \emph{one} room and asked how
many of the requests we could serve; here, every request \emph{must} be
served, and we ask how many rooms that requires. The keyword shift from
``maximum number you can fit in one resource'' to ``minimum resources needed
to fit everyone'' is the signal to switch from finish-time-only greedy
scheduling to a \emph{two-pointer sweep over both arrivals and departures}.
\end{patternbox}

\subsection*{Understanding the problem carefully}

At any instant in time, the number of platforms required is exactly the
number of trains whose interval $[\textit{arr}[i], \textit{dep}[i]]$
currently contains that instant -- i.e., the number of trains
\emph{simultaneously} present at the station. The answer to the whole
problem is simply the \textbf{maximum}, over all instants in time, of this
simultaneous count. So the task reduces to efficiently finding the peak
number of overlapping intervals.

A direct way to track this is to imagine sweeping a clock forward through
time. Every arrival event increases the number of trains currently at the
station by one; every departure event decreases it by one. If we process
all arrival and departure events in time order, the running count after
each event tells us exactly how many platforms are occupied at that moment,
and the maximum value this running count ever reaches is our answer.

\begin{ideabox}[Greedy Choice]
Sort the arrival times and the departure times \emph{independently} (we do
not need to keep them paired to the same train -- only the multiset of
arrival instants and the multiset of departure instants matter for counting
overlaps). Use two pointers, one walking through sorted arrivals and one
through sorted departures. At each step, compare the next unprocessed
arrival to the next unprocessed departure: if the arrival happens at or
before that departure, a new train has shown up while the previous one has
not yet left, so increment the platform count needed and advance the
arrival pointer; otherwise, a train has left before this conflict arises, so
decrement the platform count and advance the departure pointer. Track the
maximum platform count seen at any point.
\end{ideabox}

\subsection*{Why sorting arrivals and departures independently is valid}

It may seem surprising that we are allowed to sort $\textit{arr}$ and
$\textit{dep}$ separately, losing the information of which arrival belongs
to which train. But the quantity we are computing -- the number of
simultaneously occupied platforms at any instant -- depends only on
\emph{how many} arrival events have occurred and \emph{how many} departure
events have occurred up to that instant, not on which specific train each
event belongs to. The running difference (arrivals processed so far minus
departures processed so far) is precisely the number of trains currently at
the station, regardless of how we have paired up arrival times with
departure times internally. This is what justifies treating $\textit{arr}$
and $\textit{dep}$ as two independent sorted streams of timestamps and
sweeping through them with two pointers, taking the running difference's
maximum value.

\subsection*{Algorithm}

\begin{enumerate}
  \item Sort $\textit{arr}$ in increasing order; sort $\textit{dep}$ in
        increasing order.
  \item Initialize $i \gets 0$, $j \gets 0$, $\textit{platforms} \gets 0$,
        $\textit{maxPlatforms} \gets 0$.
  \item While $i < n$ and $j < n$:
  \begin{enumerate}
    \item If $\textit{arr}[i] \le \textit{dep}[j]$: a new train has
          arrived while the earliest still-pending train has not yet left.
          $\textit{platforms} \gets \textit{platforms}+1$;
          $\textit{maxPlatforms} \gets \max(\textit{maxPlatforms},
          \textit{platforms})$; $i \gets i+1$.
    \item Otherwise: a train has departed.
          $\textit{platforms} \gets \textit{platforms}-1$; $j \gets j+1$.
  \end{enumerate}
  \item Return \textit{maxPlatforms}.
\end{enumerate}

\begin{algorithm}[H]
\caption{Minimum Platforms}
\begin{algorithmic}[1]
\Procedure{MinPlatforms}{$\textit{arr}, \textit{dep}$}
  \State sort $\textit{arr}$ in increasing order
  \State sort $\textit{dep}$ in increasing order
  \State $i \gets 0$, $j \gets 0$, $\textit{platforms} \gets 0$, $\textit{maxPlatforms} \gets 0$
  \While{$i < n$ \textbf{and} $j < n$}
    \If{$\textit{arr}[i] \le \textit{dep}[j]$}
      \State $\textit{platforms} \gets \textit{platforms} + 1$
      \State $\textit{maxPlatforms} \gets \max(\textit{maxPlatforms}, \textit{platforms})$
      \State $i \gets i + 1$
    \Else
      \State $\textit{platforms} \gets \textit{platforms} - 1$
      \State $j \gets j + 1$
    \EndIf
  \EndWhile
  \State \Return \textit{maxPlatforms}
\EndProcedure
\end{algorithmic}
\end{algorithm}

\subsection*{Dry run}

\dryrun{Take $\textit{arr} = [900, 940, 950, 1100, 1500, 1800]$ and
$\textit{dep} = [910, 1200, 1120, 1130, 1900, 2000]$.}

Sorted: $\textit{arr} = [900, 940, 950, 1100, 1500, 1800]$ (already sorted),
$\textit{dep} = [910, 1120, 1130, 1200, 1900, 2000]$ (sorted).

\begin{itemize}
  \item $\textit{arr}[0]=900 \le \textit{dep}[0]=910$. $\textit{platforms}=1$,
        $\textit{maxPlatforms}=1$, $i=1$.
  \item $\textit{arr}[1]=940 \le \textit{dep}[0]=910$? No ($940 > 910$).
        $\textit{platforms}=0$, $j=1$.
  \item $\textit{arr}[1]=940 \le \textit{dep}[1]=1120$. $\textit{platforms}=1$,
        $\textit{maxPlatforms}=1$, $i=2$.
  \item $\textit{arr}[2]=950 \le \textit{dep}[1]=1120$. $\textit{platforms}=2$,
        $\textit{maxPlatforms}=2$, $i=3$.
  \item $\textit{arr}[3]=1100 \le \textit{dep}[1]=1120$. $\textit{platforms}=3$,
        $\textit{maxPlatforms}=3$, $i=4$.
  \item $\textit{arr}[4]=1500 \le \textit{dep}[1]=1120$? No.
        $\textit{platforms}=2$, $j=2$.
  \item $\textit{arr}[4]=1500 \le \textit{dep}[2]=1130$? No.
        $\textit{platforms}=1$, $j=3$.
  \item $\textit{arr}[4]=1500 \le \textit{dep}[3]=1200$? No.
        $\textit{platforms}=0$, $j=4$.
  \item $\textit{arr}[4]=1500 \le \textit{dep}[4]=1900$. $\textit{platforms}=1$,
        $\textit{maxPlatforms}$ stays $3$, $i=5$.
  \item $\textit{arr}[5]=1800 \le \textit{dep}[4]=1900$. $\textit{platforms}=2$,
        $\textit{maxPlatforms}$ stays $3$, $i=6$.
  \item $i = n = 6$, loop ends.
\end{itemize}

The minimum number of platforms required is $\boxed{3}$, reached around
time $1100$--$1120$, when three trains are simultaneously at the station.

\subsection*{C++ implementation}

\begin{lstlisting}
#include <bits/stdc++.h>
using namespace std;

int minPlatforms(vector<int>& arr, vector<int>& dep) {
    int n = arr.size();
    sort(arr.begin(), arr.end());
    sort(dep.begin(), dep.end());

    int i = 0, j = 0;
    int platforms = 0, maxPlatforms = 0;

    while (i < n && j < n) {
        if (arr[i] <= dep[j]) {
            platforms++;
            maxPlatforms = max(maxPlatforms, platforms);
            i++;
        } else {
            platforms--;
            j++;
        }
    }
    return maxPlatforms;
}

int main() {
    vector<int> arr = {900, 940, 950, 1100, 1500, 1800};
    vector<int> dep = {910, 1200, 1120, 1130, 1900, 2000};

    cout << "Minimum platforms required: "
         << minPlatforms(arr, dep) << endl;
    return 0;
}
\end{lstlisting}

\begin{complexitybox}
\textbf{Time Complexity.} Sorting both arrays costs $O(n \log n)$. The
two-pointer sweep afterward visits each arrival and each departure exactly
once, costing $O(n)$. The overall time complexity is
\[
  O(n \log n).
\]
\textbf{Space Complexity.} Sorting can be performed in place (or with
$O(n)$ auxiliary space depending on the sorting routine), and the sweep
itself uses only a handful of scalar variables, so the auxiliary space
complexity is $O(1)$ beyond the input arrays.
\end{complexitybox}

\begin{takeawaybox}
When a problem asks for the minimum number of identical resources needed so
that \emph{every} request can be served without waiting, think in terms of
the maximum simultaneous overlap of intervals, computed by sweeping sorted
arrival and departure timestamps with a running counter. This is a different
flavor of greedy from activity selection: there, sorting by finish time let
us discard conflicting requests one at a time; here, sorting both endpoints
lets us measure the peak load directly, with nothing discarded.
\end{takeawaybox}

\section{Job Sequencing Problem}
\label{sec:job-sequencing}

\problemstatement{We are given $n$ jobs, where the $i$-th job has a deadline
$\textit{deadline}[i]$ and a profit $\textit{profit}[i]$. Every job takes
exactly one unit of time to complete, and only one job can be processed at
any given unit of time (time slots are $1, 2, 3, \dots$). A job earns its
profit only if it is completed at or before its deadline. We must choose
which jobs to schedule, and at which time slots, so as to \textbf{maximize
total profit}; we are also asked to report how many jobs end up being
scheduled.}

\begin{patternbox}
\emph{``Each job has a deadline and a profit; schedule jobs (each taking one
unit of time) to maximize total profit''} is the signature of the
\emph{Job Sequencing Problem}. The two keywords to notice together are
\textbf{``profit, maximize''} (suggesting we should prioritize high-value
jobs first) and \textbf{``deadline''} (suggesting that, once we have decided
\emph{which} jobs to take, we still need a rule for placing them into time
slots without violating their deadlines). This combination -- rank by value,
then place greedily as late as possible -- is a different flavor of greedy
from interval scheduling, and is often paired conceptually with the Disjoint
Set Union (Union-Find) data structure for an even faster slot-finding step.
\end{patternbox}

\subsection*{Understanding the problem carefully}

Two separate questions are tangled together here: \emph{which} jobs should
we choose to do, and \emph{when} should we do each chosen job. Profit is
earned per job, independent of which legal slot it occupies, so among
several candidate sets of jobs we might choose, more total profit is always
better, with no benefit to a particular schedule beyond it being legal at
all. This suggests considering jobs in decreasing order of profit, and
greedily trying to fit each one in, since a job with higher profit is always
worth taking over a lower-profit job if there is room for only one of them.

Once we decide to take a particular job with deadline $d$, the only real
question is \emph{which} slot among $1, 2, \dots, d$ to place it in. Here a
second, equally important greedy idea applies: place the job in the
\textbf{latest} available slot not exceeding its deadline, rather than the
earliest. The reasoning is that an early slot is more valuable to keep
\emph{open}, because it can still accommodate jobs with very tight (small)
deadlines that arrive later in our profit-sorted processing order, whereas a
late slot close to $d$ is useless to any job whose deadline is smaller than
$d$ anyway.

\begin{ideabox}[Greedy Choice]
Sort all jobs in decreasing order of profit. For each job, in this order,
scan backward from its deadline toward slot $1$ and place it in the first
free slot found. If no free slot exists at or before its deadline, skip this
job permanently -- it can never be scheduled later, since slots only get
more occupied as we proceed, not less.
\end{ideabox}

\subsection*{Why processing by profit, and placing as late as possible, is correct}

Consider why we never want to revisit a profit-order decision. Once a
higher-profit job has successfully claimed a slot, no later (lower-profit)
job can ever displace it without decreasing total profit, so committing
immediately, in profit order, is safe -- this is the same "highest value
first" exchange argument used throughout greedy scheduling: if some optimal
solution skipped a higher-profit job $J$ to make room for a lower-profit job
$J'$ that conflicts with it, swapping $J$ in for $J'$ (which is always
possible, since $J$ has at least as early a deadline requirement satisfied
by any slot $J'$ could use, or an even more flexible range) cannot decrease,
and may increase, total profit.

Placing each accepted job as late as possible within its own deadline window
is the secondary greedy choice that makes the slot-assignment feasible
without any backtracking. By reserving early slots for as long as possible,
we never block a future job (processed later, because it has lower profit)
that might desperately need one of those early slots due to having a very
small deadline -- a job with deadline $1$ has nowhere else to go, so an
early slot occupied unnecessarily by an earlier, more flexible job would be
a wasted opportunity that greedy avoids by always preferring the latest
legal slot first.

\subsection*{Algorithm}

\begin{enumerate}
  \item Sort the jobs in decreasing order of profit.
  \item Let $\textit{maxDeadline}$ be the largest deadline among all jobs.
        Create a boolean array $\textit{slot}[1..\textit{maxDeadline}]$,
        all initially free.
  \item Initialize $\textit{countJobs} \gets 0$ and
        $\textit{totalProfit} \gets 0$.
  \item For each job, in decreasing profit order, with deadline $d$:
  \begin{enumerate}
    \item Scan $t$ from $\min(d, \textit{maxDeadline})$ down to $1$:
    \begin{enumerate}
      \item If $\textit{slot}[t]$ is free: mark it occupied, increment
            $\textit{countJobs}$, add this job's profit to
            $\textit{totalProfit}$, and stop scanning (break out of this
            inner loop) -- the job has been placed.
    \end{enumerate}
    \item If no free slot was found in the scan, this job is permanently
          skipped.
  \end{enumerate}
  \item Return $\textit{countJobs}$ and $\textit{totalProfit}$.
\end{enumerate}

\begin{algorithm}[H]
\caption{Job Sequencing}
\begin{algorithmic}[1]
\Procedure{JobSequencing}{\textit{jobs}}
  \State sort \textit{jobs} in decreasing order of profit
  \State $\textit{maxDeadline} \gets \max_i \textit{deadline}[i]$
  \State $\textit{slot}[1..\textit{maxDeadline}] \gets$ all free
  \State $\textit{countJobs} \gets 0$, $\textit{totalProfit} \gets 0$
  \For{each job $(d, p)$ in sorted order}
    \For{$t \gets \min(d, \textit{maxDeadline})$ \textbf{down to} $1$}
      \If{$\textit{slot}[t]$ is free}
        \State mark $\textit{slot}[t]$ occupied
        \State $\textit{countJobs} \gets \textit{countJobs} + 1$
        \State $\textit{totalProfit} \gets \textit{totalProfit} + p$
        \State \textbf{break}
      \EndIf
    \EndFor
  \EndFor
  \State \Return $(\textit{countJobs}, \textit{totalProfit})$
\EndProcedure
\end{algorithmic}
\end{algorithm}

\subsection*{Dry run}

\dryrun{Take jobs with (deadline, profit) pairs:
$(4,20), (1,10), (1,40), (1,30)$.}

Sorted by decreasing profit:
$(1,40), (1,30), (4,20), (1,10)$.

Slots $1, 2, 3, 4$ all start free.

\begin{itemize}
  \item Job $(1,40)$: scan from $t=1$ down to $1$. Slot $1$ is free; place
        it there. $\textit{countJobs}=1$, $\textit{totalProfit}=40$.
        Slots: $[1{=}\text{occupied}, 2,3,4{=}\text{free}]$.
  \item Job $(1,30)$: scan from $t=1$ down to $1$. Slot $1$ is occupied.
        No free slot found within its deadline. Skip this job.
  \item Job $(4,20)$: scan from $t=4$ down to $1$. Slot $4$ is free; place
        it there. $\textit{countJobs}=2$, $\textit{totalProfit}=60$.
        Slots: $[1{=}\text{occ}, 2,3{=}\text{free}, 4{=}\text{occ}]$.
  \item Job $(1,10)$: scan from $t=1$ down to $1$. Slot $1$ is occupied.
        Skip.
\end{itemize}

The final result is $\textit{countJobs} = 2$ jobs scheduled, with
$\textit{totalProfit} = 60$.

\subsection*{C++ implementation}

\begin{lstlisting}
#include <bits/stdc++.h>
using namespace std;

struct Job {
    int id;
    int deadline;
    int profit;
};

pair<int,int> jobSequencing(vector<Job>& jobs) {
    sort(jobs.begin(), jobs.end(), [](const Job& a, const Job& b) {
        return a.profit > b.profit;
    });

    int maxDeadline = 0;
    for (auto& job : jobs) {
        maxDeadline = max(maxDeadline, job.deadline);
    }

    vector<bool> slot(maxDeadline + 1, false); // 1-indexed
    int countJobs = 0, totalProfit = 0;

    for (auto& job : jobs) {
        for (int t = min(job.deadline, maxDeadline); t >= 1; t--) {
            if (!slot[t]) {
                slot[t] = true;
                countJobs++;
                totalProfit += job.profit;
                break;
            }
        }
    }
    return {countJobs, totalProfit};
}

int main() {
    vector<Job> jobs = {
        {1, 4, 20},
        {2, 1, 10},
        {3, 1, 40},
        {4, 1, 30}
    };

    auto result = jobSequencing(jobs);
    cout << "Jobs done: " << result.first << endl;
    cout << "Total profit: " << result.second << endl;
    return 0;
}
\end{lstlisting}

\begin{complexitybox}
\textbf{Time Complexity.} Sorting the $n$ jobs by profit costs
$O(n \log n)$. For each job, scanning backward through slots costs up to
$O(\textit{maxDeadline})$ in the worst case, giving a total of
$O(n \cdot \textit{maxDeadline})$ for the slot-placement phase. Hence the
overall time complexity is
\[
  O(n \log n + n \cdot \textit{maxDeadline}).
\]
This can be improved to $O(n \log n + n \log n) = O(n \log n)$ overall by
using a Disjoint Set Union structure (covered in Chapter~6) to jump directly
to the latest free slot at or before a given deadline, rather than scanning
slot by slot.
\textbf{Space Complexity.} We use a boolean array of size
$O(\textit{maxDeadline})$ to track slot occupancy, so the space complexity
is $O(\textit{maxDeadline})$, in addition to $O(n)$ for sorting the jobs.
\end{complexitybox}

\begin{takeawaybox}
When a problem combines \emph{``maximize total value''} with
\emph{``each item also has a deadline/capacity constraint on where it can
be placed,''} a two-layer greedy often applies: sort by value to decide
\emph{which} items to take, and within that order, place each accepted item
as late (or otherwise as flexibly) as possible so as not to block items that
have less placement freedom and are still waiting to be considered.
\end{takeawaybox}

\section{Candy}
\label{sec:candy}

\problemstatement{There are $n$ children standing in a line, and the $i$-th
child has a rating $\textit{ratings}[i]$. We must give each child at least
one candy, subject to the rule that any child with a strictly higher rating
than an \textbf{immediate neighbor} (the child directly to the left or
directly to the right) must receive strictly more candies than that
neighbor. We must find the \textbf{minimum total number of candies}
required to satisfy all of these constraints.}

\begin{patternbox}
\emph{``Each element must be strictly greater than its immediate neighbor
whenever its value is strictly greater''} -- a constraint defined purely in
terms of comparisons with the elements directly to the left and right -- is
the signature of a \emph{two-pass directional greedy}. Whenever a problem
states a rule of the form ``compare element $i$ only to $i-1$ and $i+1$, and
the rule is not symmetric (higher rating needs strictly more, but the rule
says nothing directly about how much more when ratings are equal or lower),''
the standard technique is to satisfy the left-neighbor rule and the
right-neighbor rule in two completely separate linear passes, then combine
the two results -- rather than trying to track both directions
simultaneously in a single pass, which is awkward because a child's
required candy count can be driven up by either side independently.
\end{patternbox}

\subsection*{Understanding the problem carefully}

The rule is purely \emph{local}: child $i$ needs more candies than child
$i-1$ if and only if $\textit{ratings}[i] > \textit{ratings}[i-1]$, and
child $i$ needs more candies than child $i+1$ if and only if
$\textit{ratings}[i] > \textit{ratings}[i+1]$. These are two independent
constraints, one looking left and one looking right, and a child's final
candy count must simultaneously satisfy both. This naturally suggests
solving the left-looking constraint and the right-looking constraint
separately, each with its own simple left-to-right or right-to-left greedy
sweep, and then merging the two answers by taking, for each child, whichever
requirement is stricter (i.e.\ larger).

\subsection*{The left-to-right pass}

\begin{ideabox}[Greedy Choice -- Left Pass]
Scan left to right. Give every child at least $1$ candy initially. If
$\textit{ratings}[i] > \textit{ratings}[i-1]$, this child must have strictly
more candy than the child to its left, so set $\textit{left}[i] \gets
\textit{left}[i-1] + 1$. Otherwise, the left-neighbor constraint places no
extra demand on child $i$ beyond the baseline of $1$ candy, so
$\textit{left}[i] \gets 1$.
\end{ideabox}

This single pass correctly satisfies every left-looking constraint, because
each child's requirement here only ever depends on the child immediately
before it, which has already been finalized by the time we reach the
current child -- there is no need to look further back, since any chain of
strictly increasing ratings to the left has already been correctly
propagated forward by the recursive nature of the update rule.

\subsection*{The right-to-left pass}

\begin{ideabox}[Greedy Choice -- Right Pass]
Symmetrically, scan right to left. Give every child at least $1$ candy
initially. If $\textit{ratings}[i] > \textit{ratings}[i+1]$, set
$\textit{right}[i] \gets \textit{right}[i+1] + 1$. Otherwise,
$\textit{right}[i] \gets 1$.
\end{ideabox}

\subsection*{Why taking the maximum of both passes is correct}

Once both arrays are computed, child $i$'s final candy count must be at
least $\textit{left}[i]$ (to satisfy the left-neighbor rule) and at least
$\textit{right}[i]$ (to satisfy the right-neighbor rule). Since these are
two independent lower bounds on the same quantity (the number of candies
given to child $i$), and we want the \emph{minimum} total, the smallest
value that satisfies both lower bounds simultaneously is exactly their
maximum:
\[
  \textit{candies}[i] = \max(\textit{left}[i], \textit{right}[i]).
\]
No smaller value could satisfy both constraints, and this value, taken for
every child simultaneously, does not interfere with any other child's
constraint, because each constraint here only relates two physically
adjacent children, and the relative order between $\textit{left}[i]$ and
$\textit{left}[i\pm1]$ (or $\textit{right}[i]$ and $\textit{right}[i\pm1]$)
already correctly encodes the required strict inequality wherever the
rating comparison demands it. Taking componentwise maxima of two arrays that
each individually satisfy their own one-sided constraint is guaranteed to
satisfy both constraints together, without ever needing to re-verify the
combined array.

\subsection*{Algorithm}

\begin{enumerate}
  \item Initialize arrays $\textit{left}[0..n-1]$ and
        $\textit{right}[0..n-1]$, each filled with $1$.
  \item Left pass: for $i \gets 1$ to $n-1$, if
        $\textit{ratings}[i] > \textit{ratings}[i-1]$, set
        $\textit{left}[i] \gets \textit{left}[i-1] + 1$.
  \item Right pass: for $i \gets n-2$ down to $0$, if
        $\textit{ratings}[i] > \textit{ratings}[i+1]$, set
        $\textit{right}[i] \gets \textit{right}[i+1] + 1$.
  \item Compute $\textit{total} \gets \sum_{i=0}^{n-1} \max(\textit{left}[i],
        \textit{right}[i])$.
  \item Return \textit{total}.
\end{enumerate}

\begin{algorithm}[H]
\caption{Candy Distribution}
\begin{algorithmic}[1]
\Procedure{Candy}{\textit{ratings}}
  \State $n \gets \text{length}(\textit{ratings})$
  \State $\textit{left}[0..n-1] \gets 1$, $\textit{right}[0..n-1] \gets 1$
  \For{$i \gets 1$ to $n-1$}
    \If{$\textit{ratings}[i] > \textit{ratings}[i-1]$}
      \State $\textit{left}[i] \gets \textit{left}[i-1] + 1$
    \EndIf
  \EndFor
  \For{$i \gets n-2$ \textbf{down to} $0$}
    \If{$\textit{ratings}[i] > \textit{ratings}[i+1]$}
      \State $\textit{right}[i] \gets \textit{right}[i+1] + 1$
    \EndIf
  \EndFor
  \State $\textit{total} \gets 0$
  \For{$i \gets 0$ to $n-1$}
    \State $\textit{total} \gets \textit{total} + \max(\textit{left}[i], \textit{right}[i])$
  \EndFor
  \State \Return \textit{total}
\EndProcedure
\end{algorithmic}
\end{algorithm}

\subsection*{Dry run}

\dryrun{Take $\textit{ratings} = [1, 0, 2]$.}

\textbf{Left pass.} $\textit{left} = [1, 1, 1]$ initially.
\begin{itemize}
  \item $i=1$: $\textit{ratings}[1]=0 > \textit{ratings}[0]=1$? No. Stays
        $\textit{left}[1]=1$.
  \item $i=2$: $\textit{ratings}[2]=2 > \textit{ratings}[1]=0$? Yes.
        $\textit{left}[2] = \textit{left}[1]+1 = 2$.
\end{itemize}
Result: $\textit{left} = [1, 1, 2]$.

\textbf{Right pass.} $\textit{right} = [1, 1, 1]$ initially.
\begin{itemize}
  \item $i=1$: $\textit{ratings}[1]=0 > \textit{ratings}[2]=2$? No. Stays
        $\textit{right}[1]=1$.
  \item $i=0$: $\textit{ratings}[0]=1 > \textit{ratings}[1]=0$? Yes.
        $\textit{right}[0] = \textit{right}[1]+1 = 2$.
\end{itemize}
Result: $\textit{right} = [2, 1, 1]$.

\textbf{Combine.} $\textit{candies} = [\max(1,2), \max(1,1), \max(2,1)] =
[2, 1, 2]$.

The total is $2+1+2 = 5$. This is correct: child $0$ (rating $1$) needs more
than child $1$ (rating $0$), so child $0$ gets $2$ and child $1$ gets the
baseline $1$; child $2$ (rating $2$) needs more than child $1$, so child $2$
gets $2$. No constraint links child $0$ and child $2$ directly (they are not
neighbors), so this is fully consistent and minimal.

\subsection*{C++ implementation}

\begin{lstlisting}
#include <bits/stdc++.h>
using namespace std;

int candy(vector<int>& ratings) {
    int n = ratings.size();
    vector<int> left(n, 1), right(n, 1);

    for (int i = 1; i < n; i++) {
        if (ratings[i] > ratings[i - 1]) {
            left[i] = left[i - 1] + 1;
        }
    }

    for (int i = n - 2; i >= 0; i--) {
        if (ratings[i] > ratings[i + 1]) {
            right[i] = right[i + 1] + 1;
        }
    }

    int total = 0;
    for (int i = 0; i < n; i++) {
        total += max(left[i], right[i]);
    }
    return total;
}

int main() {
    vector<int> ratings = {1, 0, 2};
    cout << "Minimum candies: " << candy(ratings) << endl;
    return 0;
}
\end{lstlisting}

\begin{complexitybox}
\textbf{Time Complexity.} We perform two linear passes (left-to-right and
right-to-left) plus one linear combination pass, each $O(n)$, giving an
overall time complexity of
\[
  O(n).
\]
\textbf{Space Complexity.} We use two auxiliary arrays, \textit{left} and
\textit{right}, each of size $n$, giving a space complexity of $O(n)$. This
can be reduced to $O(1)$ extra space with a more careful single-pass
peak-detection technique, though the two-array version is the clearest to
derive and verify correct.
\end{complexitybox}

\begin{takeawaybox}
When a constraint compares each element only to its immediate left and
right neighbors, and the comparison is directionally asymmetric (a strict
inequality matters only in one direction at a time), decompose the problem
into two independent one-directional sweeps -- one enforcing the
left-looking rule, one enforcing the right-looking rule -- and combine the
two results with a pointwise maximum (for a ``need at least this many''
constraint) or minimum (for a ``cannot exceed this many'' constraint).
\end{takeawaybox}

\section{Shortest Job First (SJF) CPU Scheduling}
\label{sec:shortest-job-first}

\problemstatement{We are given the burst times (the amount of CPU time
each process needs) of $n$ processes, all of which arrive and are ready to
run at time $0$. Only one process can run on the CPU at a time, and once we
commit to running a process we run it to completion (non-preemptive
scheduling). We must decide the order in which to run the processes so as
to \textbf{minimize the average waiting time}, where the waiting time of a
process is the amount of time it sits ready in the queue before it actually
starts executing.}

\begin{patternbox}
\emph{``Minimize average waiting time, single CPU, run-to-completion
(non-preemptive), all jobs ready at the same time''} is the textbook
\emph{Shortest Job First} scheduling rule. The keyword combination to watch
for is \textbf{``minimize the average (or total) waiting/completion
time''} together with \textbf{``one task at a time, no preemption.''}
Whenever every job is available from the very start, the order of execution
is the only choice that matters, and there is a clean, provable rule for the
best order: process the \textbf{shortest job first}. This is conceptually
close to a sorting greedy, but the insight worth absorbing is \emph{why}
ordering by burst time, rather than any other attribute, minimizes the
\emph{average} of a quantity that depends on cumulative sums.
\end{patternbox}

\subsection*{Understanding the problem carefully}

If we decide to run the processes in some order $p_1, p_2, \dots, p_n$ with
burst times $b_1, b_2, \dots, b_n$ (relabelled according to this order),
then process $p_1$ waits $0$ time units, process $p_2$ waits $b_1$ time
units (it has to wait for $p_1$ to finish), process $p_3$ waits
$b_1 + b_2$, and in general process $p_k$ waits
\[
  W_k = b_1 + b_2 + \cdots + b_{k-1}.
\]
The total waiting time is
\[
  \sum_{k=1}^{n} W_k = \sum_{k=1}^{n} \sum_{j=1}^{k-1} b_j
                      = \sum_{j=1}^{n} (n - j)\, b_j,
\]
since burst time $b_j$ is counted once for every process that comes after
it in the order -- exactly $n-j$ of them. This formula reveals the key
structural fact: \textbf{burst times that appear earlier in the order are
multiplied by a larger coefficient} ($n-1$ for the very first job, down to
$0$ for the last job). To minimize this weighted sum, we should pair the
\emph{largest} coefficients with the \emph{smallest} burst times -- which
means running the shortest jobs first.

\begin{ideabox}[Greedy Choice]
Sort the processes in increasing order of burst time, and run them strictly
in that order.
\end{ideabox}

\subsection*{Why shortest-job-first is optimal}

This is a direct instance of the \emph{rearrangement inequality}: given two
sequences of numbers, the sum of pairwise products is minimized when one
sequence is sorted increasing and the other is sorted decreasing (so that
large values are paired with small values), and maximized when both are
sorted in the same direction. Here, the coefficients $n-1, n-2, \dots, 0$
are fixed and strictly decreasing by position; to minimize
$\sum_j (n-j) b_j$, the burst times $b_j$ must be arranged in strictly
\emph{increasing} order alongside these strictly decreasing coefficients --
that is, the smallest burst time should receive the largest coefficient
(run first, waited-for by the most subsequent jobs), and the largest burst
time should receive the smallest coefficient (run last, waited-for by
nobody).

An equivalent and perhaps more intuitive exchange argument: suppose two
adjacent jobs in some schedule are out of order, i.e.\ a longer job runs
immediately before a shorter job. Swapping them leaves the waiting time of
every other job completely unchanged (the total time consumed by this pair
together is the same regardless of their internal order, so everything
after them starts at the same time either way), but it strictly decreases
the combined waiting time of \emph{this pair}, because the shorter job, now
running first, makes the longer job wait less than the longer job, when
running first, made the shorter job wait. Repeatedly fixing such
adjacent inversions -- exactly what sorting does -- can only decrease total
waiting time, so the fully sorted (shortest first) order is optimal.

\subsection*{Algorithm}

\begin{enumerate}
  \item Sort the burst times in increasing order.
  \item Initialize $\textit{waitingTime} \gets 0$ and
        $\textit{totalWait} \gets 0$.
  \item For $i \gets 0$ to $n-1$ (processing jobs in sorted order):
  \begin{enumerate}
    \item Add $\textit{waitingTime}$ to $\textit{totalWait}$ (this is the
          waiting time of the current job).
    \item $\textit{waitingTime} \gets \textit{waitingTime} +
          \textit{burst}[i]$ (the next job must wait for this one to
          finish too).
  \end{enumerate}
  \item Return $\textit{totalWait} / n$ as the average waiting time.
\end{enumerate}

\begin{algorithm}[H]
\caption{Shortest Job First Scheduling}
\begin{algorithmic}[1]
\Procedure{SJF}{\textit{burst}}
  \State sort \textit{burst} in increasing order
  \State $\textit{waitingTime} \gets 0$, $\textit{totalWait} \gets 0$
  \For{$i \gets 0$ to $n-1$}
    \State $\textit{totalWait} \gets \textit{totalWait} + \textit{waitingTime}$
    \State $\textit{waitingTime} \gets \textit{waitingTime} + \textit{burst}[i]$
  \EndFor
  \State \Return $\textit{totalWait} / n$
\EndProcedure
\end{algorithmic}
\end{algorithm}

\subsection*{Dry run}

\dryrun{Take burst times $[6, 1, 8, 9, 3]$.}

Sorted increasing: $[1, 3, 6, 8, 9]$.

\begin{itemize}
  \item $i=0$ (burst $1$): job waits $\textit{waitingTime}=0$.
        $\textit{totalWait}=0$. Update $\textit{waitingTime}=0+1=1$.
  \item $i=1$ (burst $3$): job waits $1$. $\textit{totalWait}=1$. Update
        $\textit{waitingTime}=1+3=4$.
  \item $i=2$ (burst $6$): job waits $4$. $\textit{totalWait}=5$. Update
        $\textit{waitingTime}=4+6=10$.
  \item $i=3$ (burst $8$): job waits $10$. $\textit{totalWait}=15$. Update
        $\textit{waitingTime}=10+8=18$.
  \item $i=4$ (burst $9$): job waits $18$. $\textit{totalWait}=33$. Update
        $\textit{waitingTime}=18+9=27$.
\end{itemize}

Total waiting time is $33$, and the average waiting time is
$33 / 5 = 6.6$.

\subsection*{C++ implementation}

\begin{lstlisting}
#include <bits/stdc++.h>
using namespace std;

double shortestJobFirst(vector<int>& burst) {
    sort(burst.begin(), burst.end());

    int n = burst.size();
    long long waitingTime = 0, totalWait = 0;

    for (int i = 0; i < n; i++) {
        totalWait += waitingTime;
        waitingTime += burst[i];
    }

    return (double)totalWait / n;
}

int main() {
    vector<int> burst = {6, 1, 8, 9, 3};
    cout << "Average waiting time: "
         << shortestJobFirst(burst) << endl;
    return 0;
}
\end{lstlisting}

\begin{complexitybox}
\textbf{Time Complexity.} Sorting the $n$ burst times costs
$O(n \log n)$. The subsequent single pass to accumulate waiting times costs
$O(n)$. The overall time complexity is
\[
  O(n \log n).
\]
\textbf{Space Complexity.} Sorting is performed in place, and we use only a
constant number of extra scalar variables, giving auxiliary space
complexity $O(1)$ (ignoring sort's recursion stack).
\end{complexitybox}

\begin{takeawaybox}
Whenever the objective is a sum of the form $\sum (\text{fixed
position-dependent coefficient}) \times (\text{value you get to order})$,
and the coefficients are monotonic in position, the rearrangement
inequality tells you exactly how to order the values for the minimum (or
maximum) possible sum: pair the largest coefficients with the smallest
values (or vice versa for maximization). Recognizing this pattern --
``minimize a sum of weighted, orderable quantities'' -- immediately tells
you to sort, without needing to reconstruct the exchange argument from
scratch every time.
\end{takeawaybox}

\section{Insert Interval}
\label{sec:insert-interval}

\problemstatement{We are given a list of intervals $\textit{intervals}$ that
is already sorted by start time and contains no overlaps (it represents,
say, a calendar of existing non-conflicting events). We are also given one
new interval $\textit{newInterval}$. We must insert
$\textit{newInterval}$ into the list, merging it with any existing
intervals it overlaps, so that the result is still a list of intervals
sorted by start time with no overlaps.}

\begin{patternbox}
\emph{``The existing intervals are already sorted and non-overlapping; insert
one new interval and merge where necessary''} is a strong hint that we do
not need to re-sort or re-scan everything from scratch -- the existing order
is a resource we should exploit with a single linear pass, split into three
natural phases (intervals entirely before the new one, intervals overlapping
the new one, intervals entirely after). Whenever a problem promises that the
input is already sorted, a greedy single pass that respects that order is
almost always available, replacing what would otherwise look like a
sort-and-merge operation.
\end{patternbox}

\subsection*{Understanding the problem carefully}

Because the existing intervals are sorted by start time and mutually
non-overlapping, as we scan through them from left to right, each interval
falls into exactly one of three categories relative to
$\textit{newInterval} = [\textit{ns}, \textit{ne}]$:

\begin{enumerate}
  \item \textbf{Entirely before}: the interval's end is strictly less than
        $\textit{ns}$. It cannot overlap $\textit{newInterval}$ at all, and
        because the list is sorted, every interval we have not yet visited
        could still potentially overlap, but every interval already
        confirmed to end before $\textit{ns}$ can be output immediately,
        unchanged.
  \item \textbf{Overlapping}: the interval's start is at most
        $\textit{ne}$ and its end is at least $\textit{ns}$. Any such
        interval must be merged into a single growing interval together
        with $\textit{newInterval}$ and every other overlapping interval.
  \item \textbf{Entirely after}: the interval's start is strictly greater
        than $\textit{ne}$. Once we reach such an interval, we know
        $\textit{newInterval}$ (possibly already merged with some earlier
        intervals) is finished growing, because the remaining intervals,
        being sorted, can only start even later.
\end{enumerate}

\begin{ideabox}[Greedy Choice]
Scan the existing sorted intervals once, left to right. First, copy every
interval that ends strictly before $\textit{newInterval}$ starts directly
into the answer, untouched. Next, while the current interval overlaps the
(possibly already-growing) new interval, absorb it by extending the new
interval's start to the minimum of the two starts and its end to the maximum
of the two ends. Once no more intervals overlap, insert this final merged
interval into the answer. Finally, copy every remaining interval directly
into the answer, untouched.
\end{ideabox}

\subsection*{Why a single pass suffices}

This works cleanly because of two facts guaranteed by the problem: the
existing intervals are sorted by start time, and they do not overlap each
other. Sortedness means that once we pass an interval ending before
$\textit{ns}$, no later interval can also need to be classified as
``entirely before'' incorrectly -- the boundary between ``before'' and
``possibly overlapping'' is crossed exactly once as we scan forward.
Non-overlap among the existing intervals means that, when merging the
overlapping group, we never need to worry about two existing intervals
needing to be merged with each other independently of
$\textit{newInterval}$ -- they only ever come into contact with each other
\emph{through} their shared overlap with the growing merged interval. This
is what allows the entire operation to be a single $O(n)$ pass with no
need to sort or to revisit any interval more than once.

\subsection*{Algorithm}

Write $[s_i, e_i]$ for $\textit{intervals}[i]$, and $[\textit{ns},
\textit{ne}]$ for the (possibly growing) new interval.

\begin{enumerate}
  \item Initialize an empty result list and an index $i \gets 0$.
  \item \textbf{Phase 1.} While $i < n$ and $e_i < \textit{ns}$: append
        $[s_i, e_i]$ to the result; $i \gets i+1$.
  \item \textbf{Phase 2.} While $i < n$ and $s_i \le \textit{ne}$: merge by
        setting $\textit{ns} \gets \min(\textit{ns}, s_i)$ and
        $\textit{ne} \gets \max(\textit{ne}, e_i)$; $i \gets i+1$. Once the
        loop ends, append the final merged $[\textit{ns}, \textit{ne}]$ to
        the result.
  \item \textbf{Phase 3.} While $i < n$: append $[s_i, e_i]$ to the result;
        $i \gets i+1$.
  \item Return the result.
\end{enumerate}

\begin{algorithm}[H]
\caption{Insert Interval}
\begin{algorithmic}[1]
\Procedure{InsertInterval}{\textit{intervals}, \textit{newInterval}}
  \State $\textit{result} \gets$ empty list, $i \gets 0$, $n \gets \text{length}(\textit{intervals})$
  \While{$i < n$ \textbf{and} $\textit{intervals}[i].\textit{end} < \textit{newInterval}.\textit{start}$}
    \State append $\textit{intervals}[i]$ to \textit{result}; $i \gets i+1$
  \EndWhile
  \While{$i < n$ \textbf{and} $\textit{intervals}[i].\textit{start} \le \textit{newInterval}.\textit{end}$}
    \State $\textit{newInterval}.\textit{start} \gets \min(\textit{newInterval}.\textit{start}, \textit{intervals}[i].\textit{start})$
    \State $\textit{newInterval}.\textit{end} \gets \max(\textit{newInterval}.\textit{end}, \textit{intervals}[i].\textit{end})$
    \State $i \gets i + 1$
  \EndWhile
  \State append \textit{newInterval} to \textit{result}
  \While{$i < n$}
    \State append $\textit{intervals}[i]$ to \textit{result}; $i \gets i+1$
  \EndWhile
  \State \Return \textit{result}
\EndProcedure
\end{algorithmic}
\end{algorithm}

\subsection*{Dry run}

\dryrun{Take $\textit{intervals} = [[1,3],[6,9]]$ and
$\textit{newInterval} = [2,5]$.}

\begin{itemize}
  \item \textbf{Phase 1}: check $[1,3]$. Is $3 < 2$? No. Phase 1 ends with
        nothing copied, $i=0$.
  \item \textbf{Phase 2}: check $[1,3]$. Is $1 \le 5$? Yes, overlaps. Merge:
        $\textit{newInterval} = [\min(2,1), \max(5,3)] = [1,5]$. $i=1$.
        Check $[6,9]$. Is $6 \le 5$? No. Phase 2 ends. Append $[1,5]$ to
        result.
  \item \textbf{Phase 3}: append $[6,9]$ unchanged.
\end{itemize}

The result is $[[1,5],[6,9]]$.

Now take $\textit{intervals} = [[1,2],[3,5],[6,7],[8,10],[12,16]]$ and
$\textit{newInterval} = [4,8]$.

\begin{itemize}
  \item \textbf{Phase 1}: $[1,2]$ has end $2 < 4$, copy it. $[3,5]$ has end
        $5 \not< 4$, stop. Result so far: $[[1,2]]$, $i$ points at $[3,5]$.
  \item \textbf{Phase 2}: $[3,5]$: start $3 \le 8$, overlaps. Merge:
        $\textit{newInterval}=[\min(4,3),\max(8,5)]=[3,8]$. Next, $[6,7]$:
        start $6\le8$, overlaps. Merge: $[\min(3,6),\max(8,7)]=[3,8]$
        (unchanged). Next, $[8,10]$: start $8\le8$, overlaps. Merge:
        $[\min(3,8),\max(8,10)]=[3,10]$. Next, $[12,16]$: start $12\le10$?
        No, stop. Append $[3,10]$ to result.
  \item \textbf{Phase 3}: append $[12,16]$ unchanged.
\end{itemize}

The result is $[[1,2],[3,10],[12,16]]$.

\subsection*{C++ implementation}

\begin{lstlisting}
#include <bits/stdc++.h>
using namespace std;

vector<vector<int>> insertInterval(vector<vector<int>>& intervals,
                                    vector<int>& newInterval) {
    vector<vector<int>> result;
    int n = intervals.size();
    int i = 0;

    // Phase 1: intervals ending before newInterval starts
    while (i < n && intervals[i][1] < newInterval[0]) {
        result.push_back(intervals[i]);
        i++;
    }

    // Phase 2: merge all overlapping intervals
    while (i < n && intervals[i][0] <= newInterval[1]) {
        newInterval[0] = min(newInterval[0], intervals[i][0]);
        newInterval[1] = max(newInterval[1], intervals[i][1]);
        i++;
    }
    result.push_back(newInterval);

    // Phase 3: remaining intervals
    while (i < n) {
        result.push_back(intervals[i]);
        i++;
    }

    return result;
}

int main() {
    vector<vector<int>> intervals = {{1,2},{3,5},{6,7},{8,10},{12,16}};
    vector<int> newInterval = {4, 8};

    auto result = insertInterval(intervals, newInterval);
    for (auto& iv : result) {
        cout << "[" << iv[0] << "," << iv[1] << "] ";
    }
    cout << endl;
    return 0;
}
\end{lstlisting}

\begin{complexitybox}
\textbf{Time Complexity.} Each of the $n$ existing intervals is visited
exactly once across the three phases combined (no interval is revisited),
and each visit does constant work, giving
\[
  O(n).
\]
\textbf{Space Complexity.} The result list holds at most $n+1$ intervals, so
the space complexity is $O(n)$ for the output (no other significant
auxiliary storage is used).
\end{complexitybox}

\begin{takeawaybox}
When the input is explicitly guaranteed to already be sorted, look for a
single linear pass with a small, fixed number of phases (here:
before-overlap-after) instead of reapplying a general sort-and-merge
algorithm. Exploiting a sortedness guarantee given for free in the problem
statement is itself a greedy-style optimization: it converts an
$O(n \log n)$ general merge into an $O(n)$ specialized one.
\end{takeawaybox}

\section{Merge Intervals}
\label{sec:merge-intervals}

\problemstatement{We are given an array of intervals
$\textit{intervals}[i] = [\textit{start}_i, \textit{end}_i]$, in no
particular order, and possibly overlapping with one another. We must merge
all overlapping intervals and return an array of the resulting
non-overlapping intervals that together cover exactly the same set of
points as the original collection.}

\begin{patternbox}
\emph{``Merge all overlapping intervals''} with \textbf{no} guarantee that
the input is already sorted is the unsorted counterpart of
Section~\ref{sec:insert-interval}. The keyword shift -- the absence of a
sortedness guarantee -- is the signal that we must add a sorting step
ourselves before the same kind of single linear merge pass can apply.
Whenever you see ``merge/combine all overlapping intervals'' with arbitrary
input order, the standard recipe is always: \emph{sort by start time, then
sweep once, extending or closing off a current ``running'' merged
interval}.
\end{patternbox}

\subsection*{Understanding the problem carefully}

Once the intervals are sorted by start time, two consecutive intervals (in
this sorted order) either overlap or they do not. Crucially, after sorting
by start time, if interval $j$ does not overlap the currently-growing merged
interval, no interval appearing even later (with an even larger start time)
can overlap it either, \emph{unless} merging interval $j$ itself first
extended the running interval far enough to reach that later interval --
but if $j$ did not overlap and was therefore not merged, the running
interval's end did not change, so this concern does not arise. This
guarantees that a single forward sweep, maintaining only the one
currently-open merged interval, is sufficient; we never need to revisit an
interval once we have moved past it.

\begin{ideabox}[Greedy Choice]
Sort all intervals by start time. Initialize the answer with the first
interval as the current ``running'' interval. For each subsequent interval,
in sorted order, check whether its start is at most the end of the running
interval. If so, it overlaps: extend the running interval's end to the
maximum of the two ends. If not, the running interval is finished growing:
push it into the answer, and start a new running interval from the current
one.
\end{ideabox}

\subsection*{Why sorting first makes the greedy sweep correct}

This is, in essence, the merging phase of Section~\ref{sec:insert-interval}
applied repeatedly, once for every interval, rather than just once for a
single new interval. The correctness argument is the same: because we
process intervals in increasing order of start time, the moment an interval
fails to overlap the current running interval, every future interval also
has a start time at least as large, hence also cannot reach backward to
overlap the running interval -- so the running interval can be safely
finalized and emitted. This sequential, never-look-back property is exactly
what makes a single $O(n)$ sweep (after the $O(n \log n)$ sort) sufficient,
with no need for any kind of nested loop checking all pairs of intervals
for overlap.

\subsection*{Algorithm}

\begin{enumerate}
  \item Sort $\textit{intervals}$ by start time, increasing.
  \item Initialize the result list with the first interval as the current
        running interval.
  \item For each subsequent interval $[\textit{s}, \textit{e}]$ in sorted
        order:
  \begin{enumerate}
    \item If $s \le \textit{lastEnd}$ (the end of the last interval placed
          in the result): extend that last interval's end to
          $\max(\textit{lastEnd}, e)$.
    \item Otherwise: append $[s, e]$ as a new entry in the result.
  \end{enumerate}
  \item Return the result.
\end{enumerate}

\begin{algorithm}[H]
\caption{Merge Intervals}
\begin{algorithmic}[1]
\Procedure{MergeIntervals}{\textit{intervals}}
  \State sort \textit{intervals} by start time, increasing
  \State $\textit{result} \gets [\textit{intervals}[0]]$
  \For{$i \gets 1$ to $n-1$}
    \State $\textit{last} \gets \textit{result}[\text{length}(\textit{result})-1]$
    \If{$\textit{intervals}[i].\textit{start} \le \textit{last}.\textit{end}$}
      \State $\textit{last}.\textit{end} \gets \max(\textit{last}.\textit{end}, \textit{intervals}[i].\textit{end})$
    \Else
      \State append $\textit{intervals}[i]$ to \textit{result}
    \EndIf
  \EndFor
  \State \Return \textit{result}
\EndProcedure
\end{algorithmic}
\end{algorithm}

\subsection*{Dry run}

\dryrun{Take $\textit{intervals} = [[1,3],[2,6],[8,10],[15,18]]$.}

Sorted by start time, the order is already
$[1,3],[2,6],[8,10],[15,18]$.

\begin{itemize}
  \item Initialize result $= [[1,3]]$.
  \item Process $[2,6]$: last in result is $[1,3]$, and $2 \le 3$, so
        overlap. Extend: $[1, \max(3,6)] = [1,6]$. Result $= [[1,6]]$.
  \item Process $[8,10]$: last in result is $[1,6]$, and $8 \le 6$? No.
        Append. Result $= [[1,6],[8,10]]$.
  \item Process $[15,18]$: last in result is $[8,10]$, and $15 \le 10$?
        No. Append. Result $= [[1,6],[8,10],[15,18]]$.
\end{itemize}

The final merged set is $[[1,6],[8,10],[15,18]]$.

\subsection*{C++ implementation}

\begin{lstlisting}
#include <bits/stdc++.h>
using namespace std;

vector<vector<int>> mergeIntervals(vector<vector<int>>& intervals) {
    if (intervals.empty()) return {};

    sort(intervals.begin(), intervals.end());

    vector<vector<int>> result;
    result.push_back(intervals[0]);

    for (int i = 1; i < (int)intervals.size(); i++) {
        if (intervals[i][0] <= result.back()[1]) {
            result.back()[1] = max(result.back()[1], intervals[i][1]);
        } else {
            result.push_back(intervals[i]);
        }
    }
    return result;
}

int main() {
    vector<vector<int>> intervals = {{1,3},{2,6},{8,10},{15,18}};
    auto merged = mergeIntervals(intervals);

    for (auto& iv : merged) {
        cout << "[" << iv[0] << "," << iv[1] << "] ";
    }
    cout << endl;
    return 0;
}
\end{lstlisting}

\begin{complexitybox}
\textbf{Time Complexity.} Sorting the $n$ intervals by start time costs
$O(n \log n)$. The subsequent single sweep costs $O(n)$. The overall time
complexity is
\[
  O(n \log n).
\]
\textbf{Space Complexity.} The result list holds at most $n$ intervals, so
the space complexity is $O(n)$. If in-place sorting is used and we are
allowed to modify the input, no further auxiliary space beyond the output
is required.
\end{complexitybox}

\begin{takeawaybox}
``Merge overlapping intervals'' and ``insert one interval into an
already-sorted, non-overlapping set'' (Section~\ref{sec:insert-interval})
are really the same merging logic at two different scales. Once you notice
that a single new interval inserted into a sorted list is structurally
identical to a sorted list being built up one interval at a time, both
problems collapse into the same primitive: maintain one ``currently open''
interval and only close it off when the next candidate's start exceeds its
end.
\end{takeawaybox}

\section{Non-overlapping Intervals}
\label{sec:non-overlapping-intervals}

\problemstatement{We are given an array of intervals
$\textit{intervals}[i] = [\textit{start}_i, \textit{end}_i]$. We must find
the \textbf{minimum number of intervals to remove} so that the remaining
intervals are pairwise non-overlapping (two intervals that merely touch at
an endpoint, i.e.\ one ends exactly where another begins, are \emph{not}
considered overlapping).}

\begin{patternbox}
\emph{``Minimum number of intervals to remove so that the rest are
non-overlapping''} is logically the complement of
\emph{``maximum number of non-overlapping intervals you can keep''} -- and
the latter is exactly the activity selection problem from
Section~\ref{sec:n-meetings}. Whenever a problem asks for the
\textbf{minimum number to delete/remove} to enforce some property, check
whether the property is equivalent to a maximization you already know how
to solve; if so, the answer is simply
$n - (\text{the maximum achievable count})$. Recognizing this equivalence
saves you from re-deriving a new greedy strategy from scratch.
\end{patternbox}

\subsection*{Understanding the problem carefully}

Removing the fewest intervals to make the rest non-overlapping is the same
goal, viewed backward, as keeping the most intervals while staying
non-overlapping: every interval not removed must survive, and the survivors
must be mutually non-overlapping. So if we can compute
\[
  \textit{maxKept} = \text{largest number of pairwise non-overlapping
  intervals selectable},
\]
then the answer to this problem is immediately
\[
  n - \textit{maxKept}.
\]
And computing $\textit{maxKept}$ is precisely the activity selection problem
solved in Section~\ref{sec:n-meetings}: sort by finish time, and greedily
keep an interval whenever it does not conflict with the most recently kept
one.

\begin{ideabox}[Greedy Choice]
Sort intervals by end time, increasing. Walk through them in this order,
maintaining the end time of the last interval kept. Keep the current
interval (it does not need to be removed) if its start is at least the end
of the last kept interval (recall that touching endpoints are allowed, so
the comparison is $\ge$, not strictly $>$, unlike the room-booking version
in Section~\ref{sec:n-meetings} where touching endpoints \emph{do} conflict
because a physical room cannot instantly reset). Otherwise, this interval
must be removed, since it overlaps the last interval we chose to keep.
\end{ideabox}

\subsection*{Why this reduction and greedy choice are correct}

The correctness of greedily building $\textit{maxKept}$ by always preferring
the interval that finishes earliest among the remaining candidates is
exactly the exchange argument given in Section~\ref{sec:n-meetings}: the
interval with the globally earliest finish time can always be included in
some optimal non-overlapping selection, because including it can never
block more future options than including any later-finishing alternative
would. Once that first choice is fixed, the same argument applies
recursively to the remaining intervals that start after it ends. The only
change from Section~\ref{sec:n-meetings} is the boundary condition at
matching endpoints -- this problem explicitly states that touching is not
overlapping, so the compatibility check uses
$\textit{start}_i \ge \textit{lastEnd}$ rather than a strict inequality.

\subsection*{Algorithm}

\begin{enumerate}
  \item Sort $\textit{intervals}$ by end time, increasing.
  \item Initialize $\textit{lastEnd} \gets -\infty$ (or the end of the
        first interval after selecting it) and $\textit{kept} \gets 0$.
  \item For each interval $[\textit{s}, \textit{e}]$ in sorted order:
  \begin{enumerate}
    \item If $\textit{s} \ge \textit{lastEnd}$: keep this interval.
          $\textit{kept} \gets \textit{kept}+1$;
          $\textit{lastEnd} \gets \textit{e}$.
    \item Otherwise: this interval must be removed (it overlaps the last
          kept interval); do nothing further.
  \end{enumerate}
  \item Return $n - \textit{kept}$ as the minimum number of removals.
\end{enumerate}

\begin{algorithm}[H]
\caption{Minimum Removals for Non-overlapping Intervals}
\begin{algorithmic}[1]
\Procedure{EraseOverlapIntervals}{\textit{intervals}}
  \State sort \textit{intervals} by end time, increasing
  \State $\textit{lastEnd} \gets -\infty$, $\textit{kept} \gets 0$
  \For{each $[\textit{s}, \textit{e}]$ in sorted order}
    \If{$\textit{s} \ge \textit{lastEnd}$}
      \State $\textit{kept} \gets \textit{kept} + 1$
      \State $\textit{lastEnd} \gets \textit{e}$
    \EndIf
  \EndFor
  \State \Return $n - \textit{kept}$
\EndProcedure
\end{algorithmic}
\end{algorithm}

\subsection*{Dry run}

\dryrun{Take $\textit{intervals} = [[1,2],[2,3],[3,4],[1,3]]$.}

Sorted by end time: $[1,2],[2,3],[1,3],[3,4]$ (ties between intervals ending
at the same time can be broken arbitrarily; here $[2,3]$ and $[1,3]$ both
end at $3$).

\begin{itemize}
  \item $\textit{lastEnd}=-\infty$. Interval $[1,2]$: $1 \ge -\infty$.
        Keep. $\textit{kept}=1$, $\textit{lastEnd}=2$.
  \item Interval $[2,3]$: $2 \ge 2$. Keep (touching endpoints are allowed).
        $\textit{kept}=2$, $\textit{lastEnd}=3$.
  \item Interval $[1,3]$: $1 \ge 3$? No. Remove (do not update
        \textit{lastEnd} or \textit{kept}).
  \item Interval $[3,4]$: $3 \ge 3$. Keep. $\textit{kept}=3$,
        $\textit{lastEnd}=4$.
\end{itemize}

We kept $3$ out of $4$ intervals, so the minimum number of removals is
$4 - 3 = 1$ -- specifically, removing $[1,3]$ leaves
$[1,2],[2,3],[3,4]$, which are pairwise non-overlapping (touching only at
shared endpoints).

\subsection*{C++ implementation}

\begin{lstlisting}
#include <bits/stdc++.h>
using namespace std;

int eraseOverlapIntervals(vector<vector<int>>& intervals) {
    if (intervals.empty()) return 0;

    sort(intervals.begin(), intervals.end(),
         [](const vector<int>& a, const vector<int>& b) {
             return a[1] < b[1];
         });

    int n = intervals.size();
    int kept = 1;
    int lastEnd = intervals[0][1];

    for (int i = 1; i < n; i++) {
        if (intervals[i][0] >= lastEnd) {
            kept++;
            lastEnd = intervals[i][1];
        }
    }
    return n - kept;
}

int main() {
    vector<vector<int>> intervals = {{1,2},{2,3},{3,4},{1,3}};
    cout << "Minimum removals: "
         << eraseOverlapIntervals(intervals) << endl;
    return 0;
}
\end{lstlisting}

\begin{complexitybox}
\textbf{Time Complexity.} Sorting the $n$ intervals by end time costs
$O(n \log n)$. The subsequent single pass costs $O(n)$. The overall time
complexity is
\[
  O(n \log n).
\]
\textbf{Space Complexity.} Beyond the input array (sorted in place or with
$O(n)$ auxiliary space, depending on the sort implementation), we use only
a constant number of extra scalar variables, giving auxiliary space
complexity $O(1)$.
\end{complexitybox}

\begin{takeawaybox}
Many ``minimum removals/deletions to enforce property $P$'' problems are
secretly ``maximum subset satisfying property $P$'' problems wearing a
disguise: the answer is $n$ minus the size of the largest valid subset.
Whenever you see a removal-minimization phrasing, check whether flipping it
to a keep-maximization phrasing turns it into a pattern you already
recognize -- here, it turns directly into the activity selection problem
from Section~\ref{sec:n-meetings}.
\end{takeawaybox}

